%=======================================================================
\chapter{Rate of Deformation}
%=======================================================================

Chapter 3 answered ``how much has the material been deformed''. This one
answers ``how fast is it being deformed now''. The two questions are not
the same, and the difference is the whole content of the chapter: $\bF$,
$\bC$, $\bU$ compare the present with a reference chosen once and for all,
whereas $\bD$ and $\bW$ below compare the present with the present, and
need no reference configuration at all. That is why the rate quantities
are \emph{spatial} tensors while the strain quantities were referential.

%=======================================================================
\section{Velocity gradient, stretching and spin}\label{sec:Lstretch}
%=======================================================================

Let $\bv$ be the material velocity field. By \eqr{2.19}--\eqr{2.21} and
\eqr{2.37} its referential gradient is
\begin{equation}
   \Grad\bv=\bL\bF,\qquad \bL=\grad\bv ,
   \tag{4.1}\label{eq:4.1}
\end{equation}
$\bL$ being the \emph{spatial} velocity gradient. In components, using
$\hat v_{i}=\partial\chi_{i}/\partial t$ and $F_{iA}=\chi_{i,A}$,
\begin{equation}
\begin{aligned}
   \Grad\bv&=\hat v_{i,A}\,\be_{i}\otimes\hbE_{A}
   =\frac{\partial}{\partial X_{A}}
     \Big(\frac{\partial\chi_{i}}{\partial t}\Big)\be_{i}\otimes\hbE_{A}\\
   &=\frac{\partial}{\partial t}
     \Big(\frac{\partial\chi_{i}}{\partial X_{A}}\Big)\be_{i}\otimes\hbE_{A}
   =\dot F_{iA}\,\be_{i}\otimes\hbE_{A},
\end{aligned}
   \tag{4.2}\label{eq:4.2}
\end{equation}
the interchange being permissible when $\chi_{i}(X_{B},t)$ is twice
continuously differentiable in $X_{B}$ and $t$ jointly. So
$\Grad\bv=\dot\bF$, and \eqr{4.1} gives
\begin{equation}
   \dot\bF=\bL\bF,\qquad\text{equivalently}\qquad \bL=\dot\bF\bF^{-1}.
   \tag{4.3}\label{eq:4.3}
\end{equation}

\begin{remark}[reading \eqr{4.2}: what every mark in it
means]\label{rem:hatinch4}
Four different marks appear in these three lines. Each one is there for a
reason, and the derivation stops making sense if any of them is misread.

\textbf{The caret on $\hat v_{i}$.} Remark~\ref{rem:decorations} explained
that the same velocity is three different functions, distinguished by what
they take as input. The caret selects the one that takes a reference
place, so $\hat v_{i}=\hat v_{i}(\bX,t)$. It must be that one. The
definition $\hat v_{i}=\partial\chi_{i}/\partial t$ is a partial
derivative at fixed $\bX$, and holding $\bX$ fixed is an instruction that
only makes sense in the referential description, where $\bX$ is an
independent variable. Written $\tilde v_{i}$ the same formula would say
nothing: in the spatial description the independent variable is $\bx$,
holding it fixed follows no particle at all, and $\chi_{i}$ is not a
function of $\bx$ to begin with. The caret is therefore not an ornament.
It records which of the three functions is meant, and only one of the
three makes the formula meaningful.

\textbf{The two bases.} The vectors $\be_{i}$ span the current translation
space $\Vt$, and $\hbE_{A}$ spans the referential one $\hVt$. These are
different spaces, and $\bF$ has one index in each:
$\bF=F_{iA}\,\be_{i}\otimes\hbE_{A}$. A tensor of this kind is called
two-point, because it has one foot in the reference configuration and one
in the current one. This is why the index alphabet changes, with Latin
$i,j,k$ for the present and capital $A,B,C$ for the reference. The
convention earns its place the first time a formula has to be checked. In
$\bL=\dot\bF\bF^{-1}$, for instance, $\dot\bF$ carries $\be_{i}\otimes
\hbE_{A}$ and $\bF^{-1}$ carries $\hbE_{A}\otimes\be_{j}$. The two capital
indices meet and are summed away, leaving $\be_{i}\otimes\be_{j}$. So $\bL$
has both feet in the current configuration, which is exactly the claim
\eqr{4.1} made when it called $\bL$ spatial. Had the answer come out with
a capital index left over, the formula would have been wrong, and the
alphabet is what makes that visible at a glance.

\textbf{The overdot.} $\dot F_{iA}$ is the material derivative of
Theorem~\ref{thm:matder}, meaning the derivative taken while holding the
material point fixed. When a quantity is already written as a function of
$(\bX,t)$, holding the material point fixed is the same as holding $\bX$
fixed, so the material derivative is just $\partial/\partial t$. That is
why \eqr{4.2} produces $\dot F_{iA}$ directly, with no extra term. The
convective term $(\grad\,\cdot\,)\bv$ that appears in \eqr{2.20} is the
cost of asking the same question in the spatial description, and here the
question was never asked in that description.

\textbf{The interchange of derivatives.} Swapping $\partial/\partial
X_{A}$ with $\partial/\partial t$ is the only analytic step in the
derivation. It is Clairaut's theorem: mixed partial derivatives are equal
where they are continuous. That is the entire content of the phrase
``twice continuously differentiable in $X_{B}$ and $t$ jointly''. The
hypothesis is not recited for form. Without it the two orders of
differentiation can genuinely differ, and then $\Grad\bv$ and $\dot\bF$
would be different tensors and \eqr{4.3} would fail.
\end{remark}

\begin{remark}[why $\bL$ and not $\dot\bF$]\label{rem:whyL}
$\dot\bF$ is a two-point tensor and inherits the reference configuration:
change $\kappa$ and $\dot\bF$ changes. $\bL=\dot\bF\bF^{-1}$ maps $\Vt$ to
$\Vt$ and does not: replacing $\kappa$ by another reference replaces $\bF$
by $\bF\bK$ with $\bK$ fixed in time, so
$\dot{(\bF\bK)}(\bF\bK)^{-1}=\dot\bF\bK\bK^{-1}\bF^{-1}=\bL$. The rate of
deformation is a property of the motion at the present instant alone, and
$\bL$ is the object that says so. Figure~\ref{fig:ch4comm} puts both
statements as commuting diagrams.
\end{remark}

\begin{figure}[htbp]\centering
\renewcommand{\thefigure}{A}
\begin{tikzpicture}[scale=1.0,
   ar/.style={-{Stealth[length=2.4mm]},thick},
   nd/.style={inner sep=2pt}]
% ---------- left square: Fdot = L F ----------
\node[nd] (a1) at (0,1.75)   {$\hVt$};
\node[nd] (b1) at (3.30,1.75) {$\Vt$};
\node[nd] (a2) at (0,0)      {$\hVt$};
\node[nd] (b2) at (3.30,0)   {$\Vt$};
\draw[ar] (a1)--(b1) node[midway,above,font=\small] {$\bF(t)$};
\draw[ar] (a2)--(b2) node[midway,below,font=\small] {$\bF(t+\dl t)$};
\draw[ar] (a1)--(a2) node[midway,left,font=\small]  {$\hI$};
\draw[ar] (b1)--(b2) node[midway,right,font=\small] {$\I+\dl t\,\bL$};
\node[font=\small,text=gray!75] at (1.65,0.88) {commutes};
\node[font=\small] at (1.65,-1.25) {$\Rightarrow\ \dot\bF=\bL\bF$};
% ---------- divider ----------
\draw[gray!40] (5.85,-1.70)--(5.85,2.20);
% ---------- right triangle: change of reference ----------
\node[nd] (c) at (7.05,1.75) {$\hVt_{\kappa'}$};
\node[nd] (d) at (7.05,0)    {$\hVt_{\kappa}$};
\node[nd] (e) at (10.35,0.88) {$\Vt$};
\draw[ar] (c)--(d) node[midway,left,font=\small] {$\bK$};
\draw[ar] (c)--(e) node[midway,above right=-2pt,font=\small] {$\bF\bK$};
\draw[ar] (d)--(e) node[midway,below right=-2pt,font=\small] {$\bF$};
\node[font=\footnotesize,text=gray!75] at (8.70,-0.68) {$\dot\bK=\hOb$};
\node[font=\small] at (8.70,-1.30) {$\Rightarrow\ $same $\bL$};
\end{tikzpicture}
\caption{Left: the reference space does not move, so the vertical arrow on
that side is the identity, while the configuration at $t$ is carried to the
one at $t+\dl t$ by $\bx\mapsto\bx+\dl t\,\bv$, whose gradient is
$\I+\dl t\,\bL$. Commutativity of the square is
$\bF(t+\dl t)=(\I+\dl t\,\bL)\bF(t)$, which is \eqr{4.3}. Right: changing
the reference from $\kappa$ to $\kappa'$ replaces $\bF$ by $\bF\bK$ with
$\bK$ fixed in time. The triangle commutes, and since $\dot\bK=\hOb$ the
factor cancels in $\dot\bF\bF^{-1}$: the two references disagree about
$\bF$ and about $\dot\bF$, but not about $\bL$.}
\label{fig:ch4comm}
\end{figure}

\subsection*{The polar factors and their rates}

Differentiate \eqr{3.32}, $\bF=\bR\bU=\bV\bR$, with the product rule for
tensors (Problem 1(a)):
\begin{equation}
   \dot\bF=\dot\bR\bU+\bR\dot\bU=\dot\bV\bR+\bV\dot\bR .
   \tag{4.4}\label{eq:4.4}
\end{equation}
We want $\bL$, and by \eqr{4.3} that means post-multiplying by
$\bF^{-1}$. The two forms of $\bF$ give two forms of the inverse. From
$\bF=\bR\bU$, reversing the order and inverting each factor,
$\bF^{-1}=\bU^{-1}\bR^{-1}=\bU^{-1}\bR^{t}$, the last step because $\bR$
is orthogonal so $\bR^{-1}=\bR^{t}$. From $\bF=\bV\bR$ in the same way,
$\bF^{-1}=\bR^{-1}\bV^{-1}=\bR^{t}\bV^{-1}$.

Take the first form and multiply \eqr{4.4} on the right by
$\bU^{-1}\bR^{t}$, term by term:
\[
\begin{aligned}
   \dot\bR\bU\cdot\bU^{-1}\bR^{t}
     &=\dot\bR\big(\bU\bU^{-1}\big)\bR^{t}=\dot\bR\bR^{t},\\
   \bR\dot\bU\cdot\bU^{-1}\bR^{t}
     &=\bR\dot\bU\bU^{-1}\bR^{t},
\end{aligned}
\]
where the first line used $\bU\bU^{-1}=\hI$ and the second had nothing to
cancel, since $\dot\bU$ sits between $\bR$ and $\bU^{-1}$ and tensors do
not commute. Adding the two gives the first expression below. The second
form goes the same way with $\bR^{t}\bV^{-1}$ on the right:
$\dot\bV\bR\cdot\bR^{t}\bV^{-1}=\dot\bV\bV^{-1}$ using
$\bR\bR^{t}=\I$, and $\bV\dot\bR\cdot\bR^{t}\bV^{-1}$ left as it stands.
So
\begin{equation}
   \bL=\dot\bR\bR^{t}+\bR\dot\bU\bU^{-1}\bR^{t}
      =\dot\bV\bV^{-1}+\bV\dot\bR\bR^{t}\bV^{-1}.
   \tag{4.5}\label{eq:4.5}
\end{equation}
The two expressions are equal because both are $\dot\bF\bF^{-1}$, computed
by grouping the same product in two different ways.

\begin{remark}[why we may choose the reference, and why we choose this
one]\label{rem:whykt0}
Equation \eqr{4.5} is correct but not yet usable. It expresses $\bL$
through $\bR$, $\bU$, $\bV$ together with their rates, which is six
objects where we wanted a statement about rates alone. Getting rid of the
first three is done by a manoeuvre that recurs all through the subject, so
it is worth setting out in three parts.

\emph{First, why we are free to choose at all.} Remark~\ref{rem:whyL}
established that $\bL$ does not depend on which reference configuration
was used to define $\bF$. That fact is what licenses everything that
follows. $\bL$ is a fixed object, determined by the motion at the present
instant, while $\kappa$ is a bookkeeping device we introduced ourselves.
So we may pick whichever $\kappa$ is convenient, prove a formula for
$\bL$ with it, and the formula holds for every other choice as well,
because $\bL$ was the same object all along. No generality is lost.

\emph{Second, which one to pick.} We want the terms containing $\bR$,
$\bU$, $\bV$ themselves to disappear, leaving only $\dot\bR$, $\dot\bU$,
$\dot\bV$. The deformation of a configuration relative to itself is the
identity, so the way to make $\bF$ trivial at one instant is to use the
configuration the body occupies at that very instant as the reference.
Setting $\kappa=\kappa_{t_{0}}$ gives $\bF=\shift$ at $t=t_{0}$ by
\eqr{2.79}. Uniqueness of the polar decomposition then forces
$\bR=\shift$, $\bU=\hI$ and $\bV=\I$, since those values reproduce
$\bF=\shift$ and no other values can. Every factor and every inverse in
\eqr{4.5} becomes an identity, and the expression collapses to \eqr{4.6}.

\emph{Third, the point that is easy to get wrong.} The equalities
$\bR=\shift$, $\bU=\hI$, $\bV=\I$ hold at $t_{0}$ and at no other time.
An instant later the body has moved away from $\kappa_{t_{0}}$ and the
polar factors are no longer trivial. Their derivatives are therefore not
required to vanish, and in general they do not. This is the asymmetry the
argument runs on: the values are trivial, the rates are free. It is what
allows \eqr{4.6} to isolate $\dot\bU$ and $\dot\bR$ while all the
non-differentiated factors drop out.

To see how much damage the wrong reading does, suppose one took $\bU=\hI$
to hold for all $t$. Then $\dot\bU=\hOb$, and \eqr{4.6} would give
$\bL=\dot\bR\shift^{t}$, which is skew, so $\bD=\Ob$ and every motion
would be rigid. The conclusion is absurd, and the error was a single
confusion between a value at one instant and a value on an interval.
Because the reference is chosen to sit at the present placement, this
$\kappa$ is often called the instantaneous reference, and $\bD$ and $\bW$
are said to be measured relative to the present configuration.
\end{remark}

Now choose $\kappa=\kappa_{t_{0}}$. By \eqr{2.79} and the uniqueness of the
polar factors, at $t=t_{0}$
\[
   \bR=\shift,\qquad \bU=\hI,\qquad \bV=\I ,
\]
so $\bU^{-1}=\hI$ and $\bV^{-1}=\I$, and \eqr{4.5} collapses to
\begin{equation}
   \bL=\dot\bR\shift^{t}+\shift\dot\bU\shift^{t}=\dot\bV+\dot\bR\shift^{t},
   \tag{4.6}\label{eq:4.6}
\end{equation}
all derivatives evaluated at $t_{0}$. By Problem 1(c), $\dot\bR\bR^{t}$ is
skew whenever $\bR$ is orthogonal, so $\dot\bR\shift^{t}$ is skew at
$t_{0}$; and $\shift\dot\bU\shift^{t}$ is symmetric because $\dot\bU$ is,
$\bU$ being symmetric at every instant. Splitting
\begin{equation}
   \bL=\bD+\bW;\qquad
   \bD=\Sym\bL=\tfrac12(\bL+\bL^{t}),\quad
   \bW=\Skw\bL=\tfrac12(\bL-\bL^{t}),
   \tag{4.7}\label{eq:4.7}
\end{equation}
and invoking the uniqueness of that splitting, Problem 14 of Chapter 1,
\begin{equation}
   \bD=\shift\dot\bU\shift^{t}=\dot\bV,
   \qquad
   \bW=\dot\bR\shift^{t}
   \qquad\text{at }t_{0}.
   \tag{4.8}\label{eq:4.8}
\end{equation}
So $\bD$ is the rate of stretch and $\bW$ the rate of rotation, both
measured against the configuration occupied at that instant. $\bD$ is the
\emph{stretching} tensor, $\bW$ the \emph{spin} tensor;
Figure~\ref{fig:DWsplit} shows what each half does. In components,
\begin{equation}
\begin{aligned}
   \bD&=D_{ij}\be_{i}\otimes\be_{j},\qquad
   D_{ij}=\tfrac12\big(v_{i,j}+v_{j,i}\big),\\
   \bW&=W_{ij}\be_{i}\otimes\be_{j},\qquad
   W_{ij}=\tfrac12\big(v_{i,j}-v_{j,i}\big).
\end{aligned}
   \tag{4.9}\label{eq:4.9}
\end{equation}

\begin{figure}[htbp]\centering
\renewcommand{\thefigure}{B}
\begin{tikzpicture}[scale=1.0,
   ar/.style={-{Stealth[length=2.0mm]},semithick},
   mk/.style={fill=black}]
% ============ panel 1 : D alone ============
\begin{scope}[xshift=0cm]
  \draw[guide] (0,0) circle (1.05);
  \draw[thick] (0,0) ellipse (1.42 and 0.78);
  % material marker at 45 deg on the circle -> same parameter on ellipse
  \draw[guide] (0,0)--(0.742,0.742);
  \fill (0.742,0.742) circle (1.3pt);
  \draw[thick] (0,0)--(1.004,0.552);
  \fill (1.004,0.552) circle (1.3pt);
  \draw[ar,gray!75] (0.742,0.742) to[bend right=22] (1.004,0.552);
  \node[font=\small] at (0,-1.40) {$\I+\dl t\,\bD$};
  \node[font=\footnotesize,text=gray!75,align=center,anchor=north] at (0,-1.72)
     {shape changes,\\ axes fixed};
\end{scope}
% ============ panel 2 : W alone ============
\begin{scope}[xshift=4.2cm]
  \draw[guide] (0,0) circle (1.05);
  \draw[thick] (0,0) circle (1.05);
  \draw[guide] (0,0)--(0.742,0.742);
  \fill (0.742,0.742) circle (1.3pt);
  \draw[thick] (0,0)--(0.271,1.014);
  \fill (0.271,1.014) circle (1.3pt);
  \draw[ar,gray!75] (0.742,0.742) to[bend right=22] (0.271,1.014);
  \node[font=\small] at (0,-1.40) {$\I+\dl t\,\bW$};
  \node[font=\footnotesize,text=gray!75,align=center,anchor=north] at (0,-1.72)
     {rigid turn,\\ shape untouched};
\end{scope}
% ============ panel 3 : L = D + W ============
\begin{scope}[xshift=8.4cm]
  \draw[guide] (0,0) circle (1.05);
  \draw[thick,rotate=30] (0,0) ellipse (1.42 and 0.78);
  \draw[guide] (0,0)--(0.742,0.742);
  \fill (0.742,0.742) circle (1.3pt);
  \draw[thick] (0,0)--(0.594,0.980);
  \fill (0.594,0.980) circle (1.3pt);
  \draw[ar,gray!75] (0.742,0.742) to[bend right=22] (0.594,0.980);
  \node[font=\small] at (0,-1.40) {$\I+\dl t\,\bL$};
  \node[font=\footnotesize,text=gray!75,align=center,anchor=north] at (0,-1.72)
     {both at once};
\end{scope}
\end{tikzpicture}
\caption{What the two halves of \eqr{4.7} do to a disc of material
directions at $\bx$, over a time $\dl t$, the dashed circle being the
disc as it stands now. The symmetric part $\bD$ changes its shape without
turning its axes; the skew part $\bW$ turns it without changing its shape;
and $\bL$ does both, since to first order in $\dl t$,
$\I+\dl t\,\bL=(\I+\dl t\,\bW)(\I+\dl t\,\bD)+O(\dl t^{2})$. The marked
material direction moves in all three panels, which is
Remark~\ref{rem:Wtrap}: even the shape-only panel reorients it.}
\label{fig:DWsplit}
\end{figure}

\begin{remark}[the two ``rates of stretch'' are not the same tensor]
\label{rem:trapdot}
\eqr{4.8} holds \emph{only} at the instant $t_{0}$ at which
$\kappa=\kappa_{t}$. At any other instant $\bD\neq\dot\bV$ and
$\bD\neq\shift\dot\bU\shift^{t}$; the general statements are \eqr{4.5}
together with \eqr{4.7}, and they contain $\bR$ and $\bV$ explicitly. This
is the commonest trap in the chapter: $\bD$ is \emph{not} the material
derivative of any strain measure, and in particular
$\bD\neq\dot\bE$. What is true at all times is \eqr{4.11} below, which
carries $\bF$ on both sides.
\end{remark}

\subsection*{Rates of the Cauchy--Green tensors}

Differentiate $\bC=\bF^{t}\bF$, using the product rule and Problem 1(b),
$(\bF^{t})^{\textstyle\cdot}=\dot\bF^{t}$, then \eqr{4.3}:
\begin{equation}
\begin{aligned}
   \dot\bC&=(\bF^{t}\bF)^{\textstyle\cdot}
   =\dot\bF^{t}\bF+\bF^{t}\dot\bF
   =(\bL\bF)^{t}\bF+\bF^{t}\bL\bF\\
   &=\bF^{t}\bL^{t}\bF+\bF^{t}\bL\bF
   =\bF^{t}\big(\bL+\bL^{t}\big)\bF=2\,\bF^{t}\bD\bF .
\end{aligned}
   \tag{4.10}\label{eq:4.10}
\end{equation}
The Green strain of \eqr{2.97} is $\bE=\tfrac12(\bC-\hI)$. Differentiating
it costs nothing extra, because $\hI$ is a fixed tensor and $\tfrac12$ a
fixed number, so $\dot\bE=\tfrac12\dot\bC$. Therefore
\begin{equation}
   \bF^{t}\bD\bF=\tfrac12\dot\bC=\dot\bE .
   \tag{4.11}\label{eq:4.11}
\end{equation}

\begin{remark}[why $\bC$ is the right thing to differentiate, and what
\eqr{4.11} is saying]\label{rem:whyC}
Two questions are worth asking here. Why differentiate $\bC$ rather than
$\bF$ or $\bU$, and what kind of statement is \eqr{4.11}?

\emph{Why $\bC$.} Differentiating $\bF$ gives $\dot\bF=\bL\bF$, which is
\eqr{4.3} and tells us nothing new. Differentiating $\bU$ is what one
would like, since $\bU$ is the pure stretch, but $\bU$ is defined through
a square root, $\bU=\sqrt{\bC}$, and the derivative of a tensor square
root has no simple closed form. $\bC=\bF^{t}\bF$ has neither problem. It
is a polynomial in $\bF$, so the product rule applies directly, and it is
already free of the rotation: substituting $\bF=\bR\bU$ gives
$\bC=\bU\bR^{t}\bR\bU=\bU^{2}$, with $\bR$ cancelling. So $\bC$ carries
exactly the stretch information, in a form that can be differentiated in
one line. That combination is why every rate calculation in finite
elasticity goes through $\bC$ rather than $\bU$.

\emph{What \eqr{4.11} says.} On the right sits $\dot\bE$, a referential
tensor: it lives on $\hVt$ and is what an observer using the reference
configuration would call the strain rate. On the left sits $\bD$,
a spatial tensor, sandwiched between $\bF^{t}$ and $\bF$. That sandwich is
the standard operation for moving a spatial tensor back to the reference
configuration, and it is forced by index bookkeeping alone:
$\bF^{t}$ carries $\hbE_{A}\otimes\be_{i}$, $\bD$ carries
$\be_{i}\otimes\be_{j}$, $\bF$ carries $\be_{j}\otimes\hbE_{B}$, and the
Latin indices contract in pairs to leave $\hbE_{A}\otimes\hbE_{B}$, which
is what a referential tensor must carry.

So \eqr{4.11} is not a formula relating two different physical quantities.
It says that $\bD$ and $\dot\bE$ are the same measurement expressed in two
configurations, with $\bF$ doing the translating. This also settles the
warning of Remark~\ref{rem:trapdot} in a sharper form. It is tempting to
read \eqr{4.11} as $\bD=\dot\bE$ with the $\bF$'s ignored, but the $\bF$'s
cannot be ignored: they are the only reason the two sides live in the same
space. At the one instant when $\kappa=\kappa_{t_{0}}$ we have
$\bF=\shift$, the sandwich disappears, and $\bD=\dot\bE$ holds. At every
other instant it does not.
\end{remark}
The same computation on $\bB=\bF\bF^{t}$ gives
\begin{equation}
   \dot\bB=\dot\bF\bF^{t}+\bF\dot\bF^{t}
   =\bL\bF\bF^{t}+\bF\bF^{t}\bL^{t}=\bL\bB+\bB\bL^{t}.
   \tag{4.12}\label{eq:4.12}
\end{equation}
Evaluating \eqr{4.11} at $t_{0}$, where $\bF=\shift$,
\begin{equation}
   \bD=\shift\dot\bE\shift^{t},
   \tag{4.13}\label{eq:4.13}
\end{equation}
so $\bD$ might equally be called the \emph{straining} tensor.

\begin{remark}[$\bD=\Ob$ characterises rigidity]\label{rem:Drigid}
In a rigid motion $\bF=\bQ(t)$, so $\bC=\hI$ and $\bB=\I$, both constant;
then $\dot\bC=\hOb$ and \eqr{4.11} gives $\bF^{t}\bD\bF=\hOb$, whence
$\bD=\Ob$ since $\bF$ is invertible. The converse is also true but is not
a one-line matter: $\bD=\Ob$ says only that $\bL$ is skew at each point,
and one must use that $\bL$ is a \emph{gradient} to conclude that it is
also spatially uniform. Granting that, $\bD=\Ob$ is necessary and
sufficient for rigidity, the rate analogue of \eqr{3.122}.
\end{remark}

%=======================================================================
\section{Material curves}\label{sec:matcurves}
%=======================================================================

The object to differentiate is \eqr{2.67}, $\mu\bm=\bF\bM$. It is worth
being clear about what each symbol in it is before touching it, because
the whole derivation turns on which of them depend on time.

$\bM$ is a unit vector in the reference configuration, pointing along a
chosen material fibre as that fibre was laid out in $\kappa$. It is a
label for the fibre, in the way $\bX$ is a label for a particle, and
labels do not change: $\dot\bM=\bze$ by Remark~\ref{rem:matvec}. The pair
$(\mu,\bm)$ on the left is what that fibre has become. $\bm$ is a unit
vector in the current configuration giving the fibre's present direction,
and $\mu>0$ is its present stretch, the factor by which its length has
grown. Splitting $\bF\bM$ into a length and a direction this way is a
deliberate choice, and it is the choice that makes the section work: it
separates the question ``how much longer is the fibre'' from the question
``which way is it now pointing'', and the two answers turn out to be
carried by $\bD$ and by the whole of $\bL$ respectively.

Differentiate, using the product rule on the left and
$\dot\bM=\bze$ with $\dot\bF=\bL\bF$ from \eqr{4.3} on the right:
\begin{equation}
   \dot\mu\,\bm+\mu\dot\bm=\dot\bF\bM=\bL\bF\bM=\mu\bL\bm ,
   \tag{4.14}\label{eq:4.14}
\end{equation}
the last step replacing $\bF\bM$ by $\mu\bm$, which is the equation we
started from.

Equation \eqr{4.14} contains both unknowns, $\dot\mu$ and $\dot\bm$, and
one equation cannot give both. The way to separate them is to project.
Take the inner product of \eqr{4.14} with $\bm$:
\[
   \dot\mu\,\ip\bm\bm+\mu\,\ip\bm{\dot\bm}=\mu\,\ip\bm{\bL\bm}.
\]
Two of these pieces vanish, each for its own reason.

\emph{First}, $\ip\bm\bm=1$ because $\bm$ is a unit vector, so the
coefficient of $\dot\mu$ is simply $1$.

\emph{Second}, $\ip\bm{\dot\bm}=0$. To see it, differentiate the identity
$\ip\bm\bm=1$ with respect to time. The inner product is bilinear, so it
obeys a product rule,
\[
   \frac{\dl}{\dl t}\ip\bm\bm
   =\ip{\dot\bm}\bm+\ip\bm{\dot\bm}
   =2\,\ip\bm{\dot\bm},
\]
the two terms being equal by symmetry of the inner product. The left side
is the derivative of the constant $1$, hence zero, so
$2\ip\bm{\dot\bm}=0$ and therefore $\ip\bm{\dot\bm}=0$. Nothing about the
motion was used; any unit vector, however it moves, has its derivative
perpendicular to itself.

What survives is
\[
   \dot\mu=\mu\,\ip\bm{\bL\bm}.
\]
The right side can be simplified once more. Split $\bL=\bD+\bW$ by
\eqr{4.7}. Then $\ip\bm{\bL\bm}=\ip\bm{\bD\bm}+\ip\bm{\bW\bm}$, and the
second term is zero because $\bW$ is skew: by Problem 6(a) of Chapter 1,
$\ip\bu{\bW\bu}=\ip{\bW^{t}\bu}\bu=-\ip{\bW\bu}\bu=-\ip\bu{\bW\bu}$ for
every $\bu$, and a number equal to its own negative is zero. So
\[
   \dot\mu=\mu\,\ip\bm{\bD\bm}.
\]

This already says everything, but it is customarily written in a tidier
form. Divide by $\mu$, which is legitimate because a stretch is strictly
positive, $\mu>0$, so we are never dividing by zero:
\[
   \frac{\dot\mu}{\mu}=\ip\bm{\bD\bm}.
\]
Now recognise the left side. By the chain rule, for any positive
differentiable $\mu(t)$,
\[
   \frac{\dl}{\dl t}\ln\mu
   =\frac{1}{\mu}\,\frac{\dl\mu}{\dl t}
   =\frac{\dot\mu}{\mu},
\]
so the left side is the material derivative of $\ln\mu$ and
\begin{equation}
   (\ln\mu)^{\textstyle\cdot}=\ip{\bm}{\bD\bm}.
   \tag{4.15}\label{eq:4.15}
\end{equation}

\begin{remark}[why the logarithm is the natural variable
here]\label{rem:whylog}
The step from $\dot\mu/\mu$ to $(\ln\mu)^{\textstyle\cdot}$ is only a
rewriting, and one may ask why anyone bothers. Three reasons, and they are
the same reasons the logarithmic strain is used throughout finite
elasticity.

The first is that $\dot\mu/\mu$ is the quantity with physical meaning. A
fibre growing from $1$ to $1.01$ and a fibre growing from $100$ to $101$
have the same $\dot\mu$ and are not stretching at the same rate; the
second is barely changing. Rates of growth are relative, which is why
interest, inflation and population growth are all quoted as fractions of
the current amount, and $\dot\mu/\mu$ is exactly that fraction per unit
time.

The second is that logarithms turn the composition of stretches into
addition. Stretch a fibre by $\mu_{1}$ and then by $\mu_{2}$ and the total
stretch is the product $\mu_{1}\mu_{2}$, not the sum. Taking logarithms
makes the total $\ln\mu_{1}+\ln\mu_{2}$, so $\ln\mu$ accumulates the way
one expects a strain to accumulate, and its rate is what \eqr{4.15}
delivers.

The third is that it makes the statement integrable. Reading \eqr{4.15}
backwards,
\[
   \ln\frac{\mu(t)}{\mu(t_{0})}
   =\int_{t_{0}}^{t}\ip\bm{\bD\bm}\,\dl\tau ,
\]
so the total stretch of a fibre is recovered from the history of one
diagonal component of $\bD$ along it. Had we stopped at $\dot\mu=\mu\,
\ip\bm{\bD\bm}$ we would have had a differential equation to solve rather
than an integral to evaluate.
\end{remark}

Now return to \eqr{4.14} for the other half of the answer. Solve it for
$\dot\bm$ by isolating that term,
\[
   \mu\dot\bm=\mu\bL\bm-\dot\mu\,\bm ,
\]
divide through by $\mu>0$,
\[
   \dot\bm=\bL\bm-\frac{\dot\mu}{\mu}\,\bm ,
\]
and substitute $\dot\mu/\mu=\ip\bm{\bD\bm}$ from \eqr{4.15} together with
$\bL=\bD+\bW$:
\begin{equation}
   \dot\bm=\bD\bm-(\ip\bm{\bD\bm})\bm+\bW\bm .
   \tag{4.16}\label{eq:4.16}
\end{equation}
So one equation, \eqr{4.14}, has been made to yield two, by taking its
component along $\bm$ and what is left over. That is the standard way to
get two statements out of one vector equation, and it works here because
$\bm$ is the natural direction to resolve along.

\begin{remark}[what $\dot\bm$ is: the overdot as a direction of
change]\label{rem:mdot}
The object $\dot\bm$ deserves a closer look, because the overdot is saying
something more specific here than a general rate of change.

Recall what $\bm$ is. It is the unit vector along a material fibre as that
fibre stands at the present instant. Since $\bm$ always has length one,
following it in time traces a curve on the unit sphere, and $\dot\bm$ is
the velocity along that curve. Two consequences follow at once.

\emph{It is perpendicular to $\bm$.} This was proved above in the course
of deriving \eqr{4.15}: differentiating $\ip\bm\bm=1$ gives
$\ip\bm{\dot\bm}=0$. So $\dot\bm$ carries no information about length. It
cannot, since the length is pinned at one. What it points along is the
direction in which the fibre is swinging, and its magnitude is the angular
speed of that swing, measured in radians per unit time. The information
about length was separated off into \eqr{4.15}, which is why \eqr{4.15}
and \eqr{4.16} together contain everything \eqr{4.14} contained, no more
and no less.

\emph{The middle term is what enforces the perpendicularity.} Look at the
three terms of \eqr{4.16} in turn. The term $\bW\bm$ is perpendicular to
$\bm$ automatically, since $\ip\bm{\bW\bm}=0$ for skew $\bW$. The term
$\bD\bm$ is not, in general; its component along $\bm$ is exactly
$\ip\bm{\bD\bm}$. The remaining term, $-(\ip\bm{\bD\bm})\bm$, subtracts
precisely that component and nothing else. So the formula says: apply
$\bL$ to the fibre, then remove the part that would change its length. In
compact form,
\[
   \dot\bm=(\I-\bm\otimes\bm)\bL\bm ,
\]
where $\I-\bm\otimes\bm$ is the projection onto the plane perpendicular to
$\bm$, the complement of the projection $\Proj\bn$ of Chapter 1. To check
that the compact form agrees with \eqr{4.16}, expand it:
$(\I-\bm\otimes\bm)\bL\bm=\bL\bm-\bm(\ip\bm{\bL\bm})
=\bD\bm+\bW\bm-(\ip\bm{\bD\bm})\bm$, using $\ip\bm{\bL\bm}=\ip\bm{\bD\bm}$
once more.

\emph{Which clock the dot keeps.} As everywhere in this chapter, the
overdot follows the material. So $\dot\bm$ is the turning rate an observer
riding the particle would measure, not the rate at which the direction
field $\bm(\bx,t)$ changes at a fixed place in space. The two differ by
the convective term of \eqr{2.20}, and mistaking one for the other is the
usual way to get a wrong sign in a rotating flow. Problem 10 is built on
exactly this distinction. In the potential vortex the direction field at
any fixed place is steady in time, while every material fibre passing
through that place is turning.
\end{remark}

\begin{remark}[the trap: $\bW$ is not the spin of material fibres]
\label{rem:Wtrap}
$\bm$ is a unit vector, so $\dot\bm$ is pure reorientation. One is
tempted to read \eqr{4.16} as saying that the reorientation is $\bW\bm$.
It is not: the first two terms are generally non-zero, and they are what
makes Problem 6(d) come out with $\dot\bm\neq\bze$ even though
$\bW=\Ob$. The stretching tensor turns material fibres as well, and it
turns different fibres by different amounts. Only in the special case of
\eqr{4.18} does $\bW$ tell the whole story. Figure~\ref{fig:Wtrapfig}
shows the failure concretely, in the setting of Problem 6.
\end{remark}

\begin{figure}[htbp]\centering
\renewcommand{\thefigure}{D}
\begin{tikzpicture}[scale=1.45,
   ar/.style={-{Stealth[length=2.3mm]},thick},
   sw/.style={-{Stealth[length=2.0mm]},semithick,gray!80}]
% only the upper half is drawn; the lower half is its mirror image
\draw[guide] (1.32,0) arc (0:180:1.32);
% principal axes
\draw[ar] (0,0)--(1.32,0);   \node[font=\small,right] at (1.36,0) {$\be_{1}$};
\draw[ar] (0,0)--(0,1.32);   \node[font=\small,above] at (0,1.36) {$\be_{2}$};
\draw[ar] (0,0)--(-1.32,0);  \node[font=\small,left] at (-1.36,0) {$-\be_{1}$};
% fixed directions
\fill (1.32,0) circle (1.6pt); \fill (0,1.32) circle (1.6pt);
\fill (-1.32,0) circle (1.6pt);
% three sample fibres in the first quadrant, drifting clockwise toward e1
\foreach \a in {28,52,76}{\draw[thin,gray!55] (0,0)--(\a:1.32);}
\draw[sw] (78:1.46) arc (78:58:1.46);
\draw[sw] (50:1.46) arc (50:30:1.46);
\draw[sw] (26:1.46) arc (26:8:1.46);
% second quadrant, drifting counter-clockwise toward -e1
\foreach \a in {104,128,152}{\draw[thin,gray!55] (0,0)--(\a:1.32);}
\draw[sw] (102:1.46) arc (102:122:1.46);
\draw[sw] (130:1.46) arc (130:150:1.46);
\draw[sw] (154:1.46) arc (154:172:1.46);
% labels
\node[font=\small,align=center] at (0,-0.42)
   {$\bW=\Ob$, yet $\dot\bm\neq\bze$};
\node[font=\small] at (0,-0.94)
   {$\dot\theta=-\tfrac{3}{4}\tfrac{\dot\lambda}{\lambda}\sin2\theta$};
\node[font=\footnotesize,text=gray!80,align=center] at (0,-1.44)
   {$\bullet$ = the only directions that hold still};
\end{tikzpicture}
\caption{The trap of Remark~\ref{rem:Wtrap}, in the $(\be_{1},\be_{2})$
section of Problem 6. The spin tensor vanishes identically, so nothing is
rotating in the sense of $\bW$; nevertheless every material fibre except
four drifts, and it drifts towards $\pm\be_{1}$, the direction that is
stretching fastest. The reason is \eqr{4.16}: with $\bW=\Ob$ it reads
$\dot\bm=\bD\bm-(\bm\cdot\bD\bm)\bm$, the component of $\bD\bm$
perpendicular to $\bm$, which vanishes exactly when $\bm$ is a principal
direction of $\bD$. Stretching alone reorients material, and that is why
$\bW$ cannot be called ``the spin of the material''. Only directions in
the upper half plane are drawn; the fourth fixed direction, $-\be_{2}$,
is the mirror image of the one at the top.}
\label{fig:Wtrapfig}
\end{figure}

Now specialise. Suppose that at some instant $t_{*}$ the fibre direction
$\bm$ happens to coincide with a principal vector $\bnu$ of $\bD$, so that
\[
   \bD\bnu=\eta\bnu ,\qquad \bm\big|_{t_{*}}=\bnu ,
\]
with $\eta$ the corresponding principal value. Feed this into \eqr{4.15}.
The right side becomes
\[
   \ip\bm{\bD\bm}=\ip\bnu{\bD\bnu}=\ip\bnu{\eta\bnu}
   =\eta\,\ip\bnu\bnu=\eta ,
\]
since $\bnu$ is a unit vector. So
\begin{equation}
   \eta=(\ln\mu)^{\textstyle\cdot}\big|_{t_{*}} ,
   \tag{4.17}\label{eq:4.17}
\end{equation}
which is the promised interpretation of the principal values: they are the
logarithmic stretch rates of material curves momentarily lying along the
principal directions of $\bD$. The eigenvalues of $\bD$ are not abstract
numbers; each is a rate of stretching that could be measured by marking a
fibre and timing its growth.

The same substitution does something to \eqr{4.16} as well. At $t_{*}$,
\[
   \bD\bm=\eta\bnu
   \qquad\text{and}\qquad
   (\ip\bm{\bD\bm})\bm=\eta\,\bnu ,
\]
the second because the scalar $\ip\bm{\bD\bm}$ was just computed to be
$\eta$ and $\bm=\bnu$. The two are equal, so the first two terms of
\eqr{4.16} cancel exactly and only the spin survives:
\begin{equation}
   \dot\bm\big|_{t_{*}}=\bW\bm\big|_{t_{*}} .
   \tag{4.18}\label{eq:4.18}
\end{equation}

\begin{remark}[why \eqr{4.18} is a statement about one instant
only]\label{rem:instantonly}
It is tempting to read \eqr{4.18} as saying that principal directions turn
with the spin, full stop. It does not, and the reason is worth being
careful about.

The cancellation used $\bm=\bnu$, which holds at $t_{*}$ by assumption.
But the two sides of that assumption evolve for different reasons and at
different rates. The fibre direction $\bm$ is convected by the motion,
obeying \eqr{4.16}. The principal direction $\bnu$ is whatever direction
happens to diagonalise $\bD$ at each instant, and $\bD$ is itself changing
in time, so $\bnu$ moves for reasons that have nothing to do with any
particular fibre. There is no reason for the two to stay together, and in
general they separate immediately after $t_{*}$. Once they have,
$\bD\bm$ is no longer parallel to $\bm$, the first two terms of \eqr{4.16}
no longer cancel, and \eqr{4.18} fails.

The special motions in which a fibre stays principal for an interval
rather than an instant are exactly those in which the principal axes are
frozen in the material, and Problem 8 works out what that costs.
Figure~\ref{fig:Wtrapfig} shows the generic situation, in which
$\bW=\Ob$ yet fibres still drift, precisely because they are not principal
directions and \eqr{4.18} does not apply to them.
\end{remark}

\subsection*{Rate of change of the angle between two fibres}

Let $\bM_{1},\bM_{2}$ be two material directions. By \eqr{4.16} with
$\bD+\bW=\bL$,
\begin{equation}
   \dot\bm_{1}=\bL\bm_{1}-(\ip{\bm_{1}}{\bD\bm_{1}})\bm_{1},
   \qquad
   \dot\bm_{2}=\bL\bm_{2}-(\ip{\bm_{2}}{\bD\bm_{2}})\bm_{2}.
   \tag{4.19}\label{eq:4.19}
\end{equation}
Both fibres are unit vectors, so the angle $\theta$ between them satisfies
$\cos\theta=\ip{\bm_{1}}{\bm_{2}}$. This is the equation to differentiate,
and it is worth noticing why we differentiate the cosine rather than
$\theta$ itself: $\theta$ is not directly expressible in the fibres,
whereas its cosine is a plain inner product, and inner products obey a
product rule.

Differentiate the left side by the chain rule,
$\frac{\dl}{\dl t}\cos\theta=-(\sin\theta)\dot\theta$, and the right side
by the product rule for the inner product:
\begin{equation}
\begin{aligned}
   -(\sin\theta)\dot\theta
   &=\ip{\dot\bm_{1}}{\bm_{2}}+\ip{\bm_{1}}{\dot\bm_{2}}\\
   &=\ip{\bm_{2}}{\bL\bm_{1}}
     -(\ip{\bm_{1}}{\bD\bm_{1}})(\ip{\bm_{1}}{\bm_{2}})\\
   &\qquad+\ip{\bm_{1}}{\bL\bm_{2}}
     -(\ip{\bm_{2}}{\bD\bm_{2}})(\ip{\bm_{1}}{\bm_{2}}).
\end{aligned}
   \tag{4.20}\label{eq:4.20}
\end{equation}
The second equality substituted \eqr{4.19} for each dotted fibre and then
expanded. In the first substitution, $\ip{\dot\bm_{1}}{\bm_{2}}
=\ip{\bL\bm_{1}}{\bm_{2}}-(\ip{\bm_{1}}{\bD\bm_{1}})\ip{\bm_{1}}{\bm_{2}}$,
and the first term was rewritten as $\ip{\bm_{2}}{\bL\bm_{1}}$ by symmetry
of the inner product. The second substitution is the same with the labels
exchanged.

Now specialise to the case of interest. Suppose the two fibres are
momentarily perpendicular at an instant $t_{*}$, so that
$\ip{\bm_{1}}{\bm_{2}}=0$ and $\theta=\pi/2$. Two simplifications follow
at once. The factor $\ip{\bm_{1}}{\bm_{2}}$ multiplies both of the
subtracted terms, so both vanish. And $\sin\theta=\sin(\pi/2)=1$, so the
left side is just $-\dot\theta$. What remains is
\[
   -\dot\theta\big|_{t_{*}}
   =\ip{\bm_{2}}{\bL\bm_{1}}+\ip{\bm_{1}}{\bL\bm_{2}} .
\]
Move $\bL$ off the first fibre in the first term using the definition of
the transpose, $\ip{\bm_{2}}{\bL\bm_{1}}=\ip{\bL^{t}\bm_{2}}{\bm_{1}}
=\ip{\bm_{1}}{\bL^{t}\bm_{2}}$, so that both terms act on $\bm_{2}$:
\[
   -\dot\theta\big|_{t_{*}}
   =\ip{\bm_{1}}{\bL^{t}\bm_{2}}+\ip{\bm_{1}}{\bL\bm_{2}}
   =\ip{\bm_{1}}{\big(\bL+\bL^{t}\big)\bm_{2}}
   =2\,\ip{\bm_{1}}{\bD\bm_{2}},
\]
using linearity to combine the two and then $\bL+\bL^{t}=2\bD$ from
\eqr{4.7}. Therefore
\begin{equation}
   \dot\theta\big|_{t_{*}}
   =-2\,\ip{\bm_{1}}{\bD\bm_{2}} .
   \tag{4.21}\label{eq:4.21}
\end{equation}
The sign says that a positive off-diagonal component of $\bD$ makes the
angle decrease, which is why one speaks of the fibres closing on one
another. So the
\emph{off-diagonal} components of $\bD$ are half the rates at which
initially perpendicular fibres close on one another, while by \eqr{4.17} the
\emph{diagonal} components are logarithmic stretch rates. That is the
complete interpretation of $\bD$, and Figure~\ref{fig:Dmeaning} is the
whole of it in one picture.

\begin{remark}[why $\bW$ vanished from \eqr{4.21}, and what it would have
done]\label{rem:Wdrops}
Something happened between \eqr{4.20} and \eqr{4.21} that is easy to pass
over: the spin tensor went into the calculation and did not come out. The
algebra that removed it was the combination $\bL+\bL^{t}=2\bD$, in which
the skew parts cancel. But the cancellation is not a lucky accident of the
manipulation, and seeing why shows the difference between the two halves
of $\bL$ more clearly than anything else in the chapter.

Write out the two places $\bW$ enters. Splitting $\bL=\bD+\bW$ in
\eqr{4.20}, its contribution to $-(\sin\theta)\dot\theta$ is
\[
   \ip{\bm_{2}}{\bW\bm_{1}}+\ip{\bm_{1}}{\bW\bm_{2}} .
\]
By skewness the first equals $-\ip{\bW\bm_{2}}{\bm_{1}}
=-\ip{\bm_{1}}{\bW\bm_{2}}$, so the two are negatives of one another and
their sum is zero. That is the algebra.

The reason behind the algebra is this. The angle $\theta$ is a relative
quantity. It depends on the two fibres only through how they are placed
with respect to each other, not on where either one sits in space. The
spin tensor turns the pair bodily, moving both fibres at the same angular
rate about the same axis, so it changes where each of $\bm_{1}$ and
$\bm_{2}$ points while leaving the relation between them untouched. Any
quantity built out of the pair alone must therefore be unable to detect
$\bW$, and $\dot\theta$ is such a quantity. The skewness of $\bW$ is the
algebraic form of the sentence: a rigid turn does not deform anything.

This also explains the caption of Figure~\ref{fig:Dmeaning}, which says
the picture is drawn with $\bW=\Ob$. If spin were present, the figure
would show both arms swinging round together in addition to closing on one
another, and $\dot\theta$ would come out the same. Setting $\bW=\Ob$ puts
the pair in the frame that makes the closing look symmetric, so that each
arm appears to contribute half of it. Equation \eqr{4.21} fixes only the
sum $\dot\theta$ and says nothing about how that sum is divided between
the two fibres, because the division depends on the frame in exactly the
way $\bW$ does. Remark~\ref{rem:Wtrap} is the same warning stated for one
fibre instead of two.
\end{remark}

\begin{figure}[htbp]\centering
\renewcommand{\thefigure}{C}
\begin{tikzpicture}[scale=1.28,
   ar/.style={-{Stealth[length=2.4mm]},thick}]
% ---- the two fibres as they stand now (dashed) ----
\draw[guide] (0,0)--(1.70,0);
\draw[guide] (0,0)--(0,1.70);
\node[font=\small,text=gray!70,below] at (1.62,-0.04) {$\bm_{1}$};
\node[font=\small,text=gray!70,left]  at (-0.05,1.58) {$\bm_{2}$};
% right angle marker on the current pair
\draw[gray!65,thin] (0.19,0)--(0.19,0.19)--(0,0.19);
\fill (0,0) circle (1.1pt);
\node[font=\small,below left=-3pt] at (0,0) {$\bx$};
% ---- the same two fibres a time dt later ----
\draw[ar] (0,0)--(1.980,0.278);   % grown and tilted up by 8 deg
\draw[ar] (0,0)--(0.212,1.505);   % shrunk and tilted right by 8 deg
% angle between them
\draw[gray!75,thin] (0.901,0.127) arc (8:82:0.91);
\node[font=\small] at (0.66,0.58) {$\theta$};
% annotations
\node[font=\small,anchor=west] at (2.03,0.30) {$\mu_{1}$ grows};
\node[font=\small,anchor=west] at (0.30,1.47) {$\mu_{2}$ shrinks};
\node[font=\small,anchor=north] at (1.05,-0.34)
  {$(\ln\mu_{i})^{\textstyle\cdot}=\bm_{i}\cdot\bD\bm_{i}$
   \quad\textit{(diagonal)}};
\node[font=\small,anchor=north] at (1.05,-0.90)
  {$\dot\theta=-2\,\bm_{1}\cdot\bD\bm_{2}$
   \quad\textit{(off-diagonal)}};
\end{tikzpicture}
\caption{The entire meaning of $\bD$, read off two material fibres that
are momentarily perpendicular at $\bx$. Their lengths change at the
logarithmic rates $\bm_{i}\cdot\bD\bm_{i}$, which are the diagonal
components of $\bD$ in the basis $\{\bm_{1},\bm_{2}\}$, by \eqr{4.15}.
The right angle between them closes at the rate $2\,\bm_{1}\cdot\bD\bm_{2}$,
twice the off-diagonal component, by \eqr{4.21}. Drawn with $\bW=\Ob$, so
that the closing is shared equally; in general $\bW$ turns the pair
bodily as well, and \eqr{4.21} fixes only the sum.}
\label{fig:Dmeaning}
\end{figure}

\subsection*{Vorticity}

By Problem 16 of Chapter 1 every skew $\bW$ has an axial vector $\bw$ with
$\bW\bu=\bw\times\bu$. For the spin tensor $\bw$ is the \emph{vorticity}.

\begin{remark}[where $w_{i}=\tfrac12\eps_{ijk}W_{kj}$ comes
from]\label{rem:axialformula}
This formula is not a fresh definition. It is
Proposition~\ref{prop:axial} solved for $\bw$, and both the factor
$\tfrac12$ and the reversed index pair are forced rather than chosen. Here
is the whole derivation.

Start with why an axial vector exists at all. A skew tensor satisfies
$W_{ij}=-W_{ji}$, so its diagonal entries vanish and the entries below the
diagonal are determined by those above it. In three dimensions that leaves
three independent numbers, namely $W_{12}$, $W_{13}$, $W_{23}$. A vector
in three dimensions also has three components. The correspondence between
skew tensors and vectors rests on this coincidence and exists in no other
dimension: in four dimensions a skew tensor has six independent entries
and no vector to match them.

The correspondence is pinned down by requiring $\bW\bu=\bw\times\bu$ for
every $\bu$. Write both sides in components. The left side is $W_{ij}
u_{j}$. For the right side, the cross product in components is
$(\bw\times\bu)_{i}=\eps_{ijk}w_{j}u_{k}$, so
\[
   W_{ij}u_{j}=\eps_{ijk}w_{j}u_{k}.
\]
The free index is $i$ and $\bu$ is arbitrary, so the coefficients of
$u_{j}$ must agree on both sides. On the right, rename the summed index
$k$ to $j$ and the summed index $j$ to $k$, which is allowed because both
are dummies:
\[
   \eps_{ijk}w_{j}u_{k}=\eps_{ikj}w_{k}u_{j}.
\]
Matching coefficients of $u_{j}$ now gives
\[
   W_{ij}=\eps_{ikj}w_{k}=-\eps_{ijk}w_{k},
\]
the last step using the antisymmetry $\eps_{ikj}=-\eps_{ijk}$. This is the
statement recorded in Proposition~\ref{prop:axial}, and it expresses $\bW$
in terms of $\bw$.

To go the other way, contract both sides with $\eps_{ijl}$ and sum over
$i$ and $j$:
\[
   \eps_{ijl}W_{ij}=-\eps_{ijl}\eps_{ijk}w_{k}.
\]
The double contraction is handled by Theorem~\ref{thm:epsdelta}, which is
stated with the repeated pair in the last two slots. A cyclic shift puts
each factor into that form, $\eps_{ijl}=\eps_{lij}$ and
$\eps_{ijk}=\eps_{kij}$, so
\[
   \eps_{ijl}\eps_{ijk}=\eps_{lij}\eps_{kij}=2\dd_{lk},
\]
and therefore
\[
   \eps_{ijl}W_{ij}=-2\dd_{lk}w_{k}=-2w_{l}
   \qquad\Longrightarrow\qquad
   w_{l}=-\tfrac12\eps_{ijl}W_{ij}=-\tfrac12\eps_{lij}W_{ij}.
\]

Now the factor $\tfrac12$ has an explanation. It is the $2$ in
$\eps_{lij}\eps_{kij}=2\dd_{lk}$, and that $2$ counts the two orderings of
the summed pair of indices. Each independent entry of $\bW$ is reached
twice by the contraction, once as $W_{ij}$ and once as $W_{ji}$, and since
the two are negatives of one another and $\eps$ also changes sign, the two
contributions add rather than cancel. Dividing by $2$ undoes the double
count. Nobody chose the factor for tidiness.

Finally, the form used below is $w_{i}=\tfrac12\eps_{ijk}W_{kj}$, with the
indices of $W$ in the order $kj$ rather than $jk$. This is the same
statement written differently. Relabel the dummies $j\leftrightarrow k$ in
$\eps_{ijk}W_{kj}$ to get $\eps_{ikj}W_{jk}$, then use
$\eps_{ikj}=-\eps_{ijk}$:
\[
   \tfrac12\eps_{ijk}W_{kj}=\tfrac12\eps_{ikj}W_{jk}
   =-\tfrac12\eps_{ijk}W_{jk},
\]
which is what the boxed result above gives with $l$ renamed to $i$.
Writing $W_{kj}$ instead of $W_{jk}$ is a way of carrying the minus sign in
the index order rather than in front of the expression, and it is why
\eqr{4.22} runs through with no stray signs to keep track of.
\end{remark}

In components, with $w_{i}=\tfrac12\eps_{ijk}W_{kj}$ and \eqr{4.9},
\begin{equation}
\begin{aligned}
   w_{i}&=\tfrac12\eps_{ijk}W_{kj}
   =\tfrac14\eps_{ijk}\big(v_{k,j}-v_{j,k}\big)\\
   &=\tfrac14\big(\eps_{ijk}v_{k,j}+\eps_{ikj}v_{j,k}\big)\\
   &=\tfrac14\big(\eps_{ijk}v_{k,j}+\eps_{ijk}v_{k,j}\big)
    =\tfrac12\eps_{jki}v_{k,j},
\end{aligned}
   \tag{4.22}\label{eq:4.22}
\end{equation}
Each line deserves a word, since three different index manipulations are
used in quick succession.

The first line substituted $W_{kj}=\tfrac12(v_{k,j}-v_{j,k})$ from
\eqr{4.9}, which contributed the extra $\tfrac12$ and made the total
$\tfrac14$.

The second line rewrote the subtraction as an addition. Look at the second
piece, $-\tfrac14\eps_{ijk}v_{j,k}$. Using antisymmetry in the last two
slots, $-\eps_{ijk}=\eps_{ikj}$, so the piece equals
$+\tfrac14\eps_{ikj}v_{j,k}$. The minus sign has been absorbed into the
order of the indices on $\eps$, which is the same trick
Remark~\ref{rem:axialformula} described.

The third line renamed the dummy indices in that second piece,
$j\leftrightarrow k$, which is always permitted for summed indices and
changes nothing. It turns $\eps_{ikj}v_{j,k}$ into $\eps_{ijk}v_{k,j}$,
which is exactly the first piece. So the two pieces are identical, they
add, and $\tfrac14+\tfrac14=\tfrac12$.

The fourth line applied the cyclic symmetry $\eps_{ijk}=\eps_{jki}$ to put
the free index $i$ last, which is the form in which the curl was written
in Chapter 1. Comparing with $(1.165)$ there,
\begin{equation}
   \bw=\tfrac12\crl\bv ,
   \tag{4.23}\label{eq:4.23}
\end{equation}
$\crl$ being the spatial curl.

\begin{remark}[the factor $\tfrac12$, and what vorticity
measures]\label{rem:halfcurl}
The factor of $\tfrac12$ in \eqr{4.23} is a standing source of confusion,
because different books define vorticity as $\bw$ and as $\crl\bv$ without
always saying which. The way to keep it straight is to remember what each
one counts, and the cleanest test is rigid rotation.

Let the body rotate rigidly at angular velocity $\bOm$, so that
$\bv=\bOm\times\br$ with $\br$ measured from a point on the axis. Then
$\bL=\grad\bv$ is the skew tensor whose axial vector is $\bOm$, so
$\bD=\Ob$, $\bW=\bL$, and by definition $\bw=\bOm$. The vorticity is the
angular velocity itself, with no factor.

Now compute the curl of the same field. Since $\bw=\tfrac12\crl\bv$ and
$\bw=\bOm$, we get $\crl\bv=2\bOm$: the curl is \emph{twice} the angular
velocity. The factor is real and it is in the curl, not in the vorticity.

The reason is that $\crl\bv$ measures the circulation per unit area around
a small loop, and a rigidly rotating disc of radius $a$ has circulation
$\oint\ip\bv{\dl\bx}=(\Omega a)(2\pi a)=2\pi\Omega a^{2}$ around its rim,
while its area is $\pi a^{2}$. The ratio is $2\Omega$, not $\Omega$.
Circulation counts the whole way round; angular velocity counts the
turning rate. Defining $\bw$ with the $\tfrac12$ is what makes the
vorticity equal to the angular velocity of the local rigid rotation, which
is the physically useful normalisation and the one this book uses.

The disc of the next section is the concrete version of this argument, and
Problem 10 shows how badly intuition fails when the flow is not rigid: a
potential vortex has curved streamlines and zero vorticity everywhere off
the axis.
\end{remark}

%=======================================================================
\section{Example: a spinning disc}\label{sec:disc}
%=======================================================================

A disc spins about a fixed centre $\bx\in\kappa_{t}$. With $\by$ a point
of the rim and $\br=\by-\bx$,
\begin{equation}
   \dot\br=\omega\,\bn\times\br ,
   \tag{4.24}\label{eq:4.24}
\end{equation}
$\omega$ the constant rotation rate and $\bn$ a fixed unit normal to the
disc (Figure~\ref{fig:disc}). Taking $\bv(\bx,t)=\bze$,
\begin{equation}
   \bv(\by,t)-\bv(\bx,t)=\omega\,\bn\times(\by-\bx).
   \tag{4.25}\label{eq:4.25}
\end{equation}
The right side is linear in $\by-\bx$, so there is a spatially uniform
$\bL(t)$ with $\bv(\by,t)-\bv(\bx,t)=\bL(t)(\by-\bx)$. It is skew, so
$\bL=\bW$, $\bD=\Ob$, and the axial vector is
\begin{equation}
   \bw(t)=\omega\,\bn .
   \tag{4.26}\label{eq:4.26}
\end{equation}

\begin{figure}[htbp]\centering
\bookfig{1}
\begin{tikzpicture}[scale=1.15,
  vec/.style={-{Stealth[length=2.4mm]},thick},
  lbl/.style={fill=white,inner sep=1.2pt,font=\small}]
\draw[thick] (0,0) ellipse (2.05 and 0.72);
\fill[gray!10] (0,0) ellipse (2.05 and 0.72);
\draw[thick] (0,0) ellipse (2.05 and 0.72);
\fill (0,0) circle (1.0pt);
\node[lbl,below left=0pt] at (0,0) {$\bx$};
\coordinate (Y) at (1.62,0.44);
\fill (Y) circle (1.0pt);
\node[lbl,below right=-1pt and 0pt] at (Y) {$\by$};
\draw[vec] (0,0)--(Y);
\node[lbl] at (0.84,0.36) {$\br$};
% the normal
\draw[vec] (0,0)--(0,1.75);
\node[lbl,above] at (0,1.78) {$\bn$};
% the velocity at y: omega n x r, tangent to the rim
\draw[vec] (Y)--++(-0.30,0.86);
\node[lbl] at (1.30,1.42) {$\omega\,\bn\times\br$};
% sense of rotation
\draw[-{Stealth[length=2.2mm]},gray!70,thick]
   (-1.35,0.62) arc (150:105:2.0 and 0.95);
\node[lbl,text=gray!85] at (-1.05,1.02) {$\omega$};
\end{tikzpicture}
\caption{The spinning disc. Every point of the body moves with
$\bv(\by)-\bv(\bx)=\omega\bn\times(\by-\bx)$, which is linear in
$\by-\bx$ with a skew coefficient. Hence $\bD=\Ob$: the motion is rigid.
The vorticity is $\bw=\omega\bn$, the angular velocity itself, because
$\crl\bv$ of a rigid field is twice the angular velocity and \eqr{4.23}
halves it.}
\label{fig:disc}
\end{figure}

In components \eqr{4.25} reads
\begin{equation}
   v_{i}(\by,t)-v_{i}(\bx,t)=\omega\,\eps_{ijk}n_{j}(y_{k}-x_{k}),
   \tag{4.27}\label{eq:4.27}
\end{equation}
so, matching against $L_{ik}(y_{k}-x_{k})$,
\[
   L_{ik}=\omega\,\eps_{ijk}n_{j}.
\]
Then $L_{ki}=\omega\eps_{kji}n_{j}=-\omega\eps_{ijk}n_{j}=-L_{ik}$, so
$D_{ik}=0$ and $W_{ik}=\omega\eps_{ijk}n_{j}$: the motion is rigid. The
vorticity follows from \eqr{4.22},
\begin{equation}
\begin{aligned}
   \bw&=w_{j}\be_{j}=\tfrac12\eps_{jki}v_{i,k}\,\be_{j}
   =\tfrac12\eps_{jki}L_{ik}\,\be_{j}\\
   &=\tfrac12\,\omega\,\eps_{jki}\eps_{ilk}n_{l}\,\be_{j}
   =\tfrac12\,\omega\,(2\dd_{jl})\,n_{l}\,\be_{j}
   =\omega\,n_{j}\be_{j}=\omega\bn ,
\end{aligned}
   \tag{4.28}\label{eq:4.28}
\end{equation}
the contraction being $\eps_{jki}\eps_{ilk}=\eps_{ijk}\eps_{ilk}
=2\dd_{jl}$ after cyclic shifts. This confirms \eqr{4.26}.

\begin{remark}[two misprints in the printed \eqr{4.27}--\eqr{4.28}]
\label{rem:ch4misprint}
The text writes ``\,$L_{ij}=\omega e_{ijk}n_{j}$\,'' and, two lines later,
``\,$W_{ij}=\omega e_{ijk}n_{j}$\,''. Neither can be right as printed:
the index $j$ occurs three times, once free on the left and twice on the
right, so the equation is not index balanced. Matching \eqr{4.27} against
$v_{i}-v_{i}=L_{ik}(y_{k}-x_{k})$ shows that the free indices are $i$ and
$k$ and that $j$ is the summed one, giving
$L_{ik}=\omega\eps_{ijk}n_{j}$ and $W_{ik}=\omega\eps_{ijk}n_{j}$, as
used above. Nothing downstream changes.

Second, the printed chain for \eqr{4.28} carries $\omega$ in the line
$\tfrac12\omega\eps_{jki}L_{ik}\be_{j}$ and again in the following line.
Since $L$ already contains $\omega$, the factor enters exactly once, at
the substitution of $L$, as written above. Again the final answer
$\bw=\omega\bn$ is unaffected.
\end{remark}

%=======================================================================
\section{Rigid-body motions}\label{sec:rigidrate}
%=======================================================================

A rigid motion is \eqr{3.121}, $\bx=\bQ(t)\bX+\bc(t)$, with $\bQ$
orthogonal. Then $\bF=\bQ$ and, by \eqr{4.3} with $\bQ^{-1}=\bQ^{t}$,
\begin{equation}
   \bL=\dot\bQ\bQ^{t}.
   \tag{4.29}\label{eq:4.29}
\end{equation}
By Problem 1(c) this is skew, so
\[
   \bD=\Ob,\qquad \bW=\dot\bQ\bQ^{t},
\]
and $\bW$ is spatially uniform, depending on $t$ alone. To recover the
velocity field, integrate at fixed $t$; uniformity of $\bW$ is what lets
it pass inside the differential, exactly as at \eqr{3.99}:
\begin{equation}
   \dl\bv=\bL\,\dl\bx=\bW(t)\,\dl\bx=\dl\big(\bW\bx\big),
   \tag{4.30}\label{eq:4.30}
\end{equation}
so $\dl(\bv-\bW\bx)=\bze$ and the bracket is independent of $\bx$:
\begin{equation}
   \bv(\bx,t)=\bW(t)\bx+\bd(t)=\bw(t)\times\bx+\bd(t),
   \tag{4.31}\label{eq:4.31}
\end{equation}
with $\bd(t)\in\Vt$ spatially uniform. Subtracting the same relation at
two points,
\begin{equation}
   \bv(\by,t)-\bv(\bx,t)=\bw(t)\times(\by-\bx),
   \tag{4.32}\label{eq:4.32}
\end{equation}
of which the spinning disc is the special case $\bw=\omega\bn$.

Alternatively differentiate \eqr{3.121} directly and eliminate $\bX$ by
$\bX=\bQ^{t}(\bx-\bc)$:
\begin{equation}
   \bv=\dot\bQ\bX+\dot\bc=\dot\bQ\bQ^{t}(\bx-\bc)+\dot\bc
     =\bW\bx+\big(\dot\bc-\bW\bc\big),
   \tag{4.33}\label{eq:4.33}
\end{equation}
which is \eqr{4.31} with $\bd=\dot\bc-\bW\bc$.

\begin{remark}[what survives from Chapter 3]\label{rem:ch4summary}
Three statements now say the same thing about rigidity, at three levels:
\[
\begin{aligned}
   \text{geometry:}&\quad |\bchi(\bY,t)-\bchi(\bX,t)|=|\bY-\bX|,
     &&\eqr{3.113}\\
   \text{deformation:}&\quad \bU=\hI,\ \ \bC=\hI,\ \ \bE=\hOb,
     &&\eqr{3.122}\\
   \text{rate:}&\quad \bD=\Ob .&&
\end{aligned}
\]
The third is the useful one in practice, because it is a condition on the
present velocity field alone and needs no reference configuration. It is
also the reason constitutive laws for fluids are written in $\bD$: a fluid
cannot feel how far it has come from a reference state, only how fast it
is being distorted now.
\end{remark}

%=======================================================================
\setcounter{problemenv}{0}
\section{Problems}\label{sec:problems:c4}
%=======================================================================

%-----------------------------------------------------------------------
\begin{problem}[Rules for time derivatives of tensors]
(a) Product rule. (b) Differentiation commutes with transposition.
(c) If $\bR$ is orthogonal, $\dot\bR\bR^{t}$ is skew. (d) Differentiation
does \emph{not} commute with inversion; find
$(\bA^{-1})^{\textstyle\cdot}$.
\end{problem}

\begin{solution}
Every part of this problem is settled the same way. A tensor equation is
nine scalar equations in components, the components are ordinary functions
of $t$, and the ordinary rules of calculus apply to them. The only care
needed is that tensors do not commute, so the \emph{order} of factors must
be preserved when the component statement is read back as a tensor
statement. All four parts are used repeatedly in this chapter and the
next, and (a) and (c) in particular are what \eqr{4.4} and \eqr{4.6} rest
on.

\textbf{(a)} In components $(\bA\bB)_{ij}=A_{ik}B_{kj}$, a sum of products
of scalar functions of $t$, so the scalar product rule applies to each
term of the sum:
\[
   \big[(\bA\bB)^{\textstyle\cdot}\big]_{ij}
   =\big(A_{ik}B_{kj}\big)^{\textstyle\cdot}
   =\dot A_{ik}B_{kj}+A_{ik}\dot B_{kj}.
\]
Read the two terms back as tensors. In the first, the free indices $i$ and
$j$ sit on $\dot A$ and $B$ in that order with $k$ summed between them, so
it is $\dot\bA\bB$; the second is $\bA\dot\bB$ by the same reading. Hence
\[
   (\bA\bB)^{\textstyle\cdot}=\dot\bA\bB+\bA\dot\bB .
\]
The order matters: $\bB\dot\bA$ would be a different tensor, and writing
the rule with the factors commuted is the commonest slip in the subject.

\textbf{(b)} $(\bA^{t})_{ij}=A_{ji}$, so
$\big[(\bA^{t})^{\textstyle\cdot}\big]_{ij}=\dot A_{ji}
=(\dot\bA^{t})_{ij}$.

\textbf{(c)} Differentiate $\bR\bR^{t}=\I$ by (a) and (b):
\[
   \dot\bR\bR^{t}+\bR\dot\bR^{t}=\Ob
   \ \Longrightarrow\
   \dot\bR\bR^{t}=-\bR\dot\bR^{t}=-\big(\dot\bR\bR^{t}\big)^{t},
\]
which is skewness.

\textbf{(d)} Differentiate $\bA\bA^{-1}=\I$:
\[
   \dot\bA\bA^{-1}+\bA\big(\bA^{-1}\big)^{\textstyle\cdot}=\Ob
   \ \Longrightarrow\
   \boxed{\ \big(\bA^{-1}\big)^{\textstyle\cdot}
   =-\bA^{-1}\dot\bA\bA^{-1}.\ }
\]
This is not $(\dot\bA)^{-1}$: the two differ already for
$\bA=t\I$, where $(\bA^{-1})^{\textstyle\cdot}=-t^{-2}\I$ while
$(\dot\bA)^{-1}=\I$. The scalar case $\dl(1/a)/\dl t=-\dot a/a^{2}$ is the
same statement, with the two copies of $\bA^{-1}$ collapsing to $a^{-2}$
only because scalars commute.
\end{solution}

%-----------------------------------------------------------------------
\begin{problem}[$\dot J=J\tr\bL$, and $\partial J/\partial F_{iA}=F^{*}_{iA}$]
(a) Differentiate $J\tbox\bu\bv\bw=\tbox{\bF\bu}{\bF\bv}{\bF\bw}$ at fixed
$\bu,\bv,\bw$ and use the box-product definition of the trace to get
$\dot J/J=\tr\bL$.
(b) Compare with $\dot J=(\partial J/\partial F_{iA})\dot F_{iA}$ to
conclude $\partial J/\partial F_{iA}=F^{*}_{iA}$.
\end{problem}

\begin{solution}
\textbf{(a)} The box product is trilinear, so its derivative obeys a
three-term product rule. With $\bu,\bv,\bw$ fixed and $\dot\bF=\bL\bF$,
\[
\begin{aligned}
   \dot J\,\tbox\bu\bv\bw
   &=\tbox{\dot\bF\bu}{\bF\bv}{\bF\bw}
    +\tbox{\bF\bu}{\dot\bF\bv}{\bF\bw}
    +\tbox{\bF\bu}{\bF\bv}{\dot\bF\bw}\\
   &=\tbox{\bL\ba}{\bb}{\bc}+\tbox{\ba}{\bL\bb}{\bc}
    +\tbox{\ba}{\bb}{\bL\bc},
\end{aligned}
\]
where $\ba=\bF\bu$, $\bb=\bF\bv$, $\bc=\bF\bw$. The last line is exactly
Definition~\ref{def:inv3}$_{1}$ of the trace, so it equals
$(\tr\bL)\tbox\ba\bb\bc=(\tr\bL)J\tbox\bu\bv\bw$. Cancelling the non-zero
box product,
\[
   \boxed{\ \dot J=J\,\tr\bL=J\,\tr\bD\ }
\]
the second form because $\tr\bW=0$ for skew $\bW$. So $J$ is constant in
time exactly when $\dvg\bv=\tr\bL=0$: that is what \emph{isochoric} means
for a motion.

\textbf{(b)} By the chain rule, $J$ being a function of the nine $F_{iA}$,
\[
   \dot J=\frac{\partial J}{\partial F_{iA}}\,\dot F_{iA}.
\]
On the other hand, by (a) and $L_{ii}=\dot F_{iA}F^{-1}_{Ai}$ from
\eqr{4.3},
\[
   \dot J=J\,L_{ii}=J\,F^{-1}_{Ai}\,\dot F_{iA}
   =\big(J\,F^{-t}_{iA}\big)\dot F_{iA}.
\]
The two hold for every $\dot F_{iA}$, which are arbitrary, so the
coefficients agree:
\[
   \frac{\partial J}{\partial F_{iA}}=J\,F^{-t}_{iA}=F^{*}_{iA},
\]
the last step by \eqr{2.55}. So the cofactor is the derivative of the
determinant, which is the identity that produces the Piola stress in
Chapter 6.
\end{solution}

%-----------------------------------------------------------------------
\begin{problem}[A material derivative from a spatial field]
With $\theta(\bx,t)=(A\cos t)/r$, $r^{2}=x_{1}^{2}+x_{2}^{2}$, and
$v_{1}=x_{1}^{2}x_{2}+x_{3}^{2}$, $v_{2}=-x_{1}^{3}-x_{1}x_{2}^{2}$,
$v_{3}=0$, find $\dot\theta$.
\end{problem}

\begin{solution}
Use \eqr{2.17}, $\dot\theta=\theta'+\ip{\grad\theta}{\bv}$.

\emph{Spatial time derivative}, at fixed $\bx$:
\[
   \theta'=\frac{\partial}{\partial t}\Big(\frac{A\cos t}{r}\Big)
   =-\frac{A\sin t}{r}.
\]

\emph{Spatial gradient.} From $r^{2}=x_{1}^{2}+x_{2}^{2}$,
$r_{,1}=x_{1}/r$, $r_{,2}=x_{2}/r$, $r_{,3}=0$. Then
$\theta=A\cos t\,r^{-1}$ gives
\[
   \theta_{,i}=-A\cos t\,r^{-2}r_{,i}
   \ \Longrightarrow\
   \grad\theta=-\frac{A\cos t}{r^{3}}\big(x_{1}\be_{1}+x_{2}\be_{2}\big).
\]

\emph{Convective term.} Only $v_{1},v_{2}$ contribute:
\[
\begin{aligned}
   x_{1}v_{1}+x_{2}v_{2}
   &=x_{1}\big(x_{1}^{2}x_{2}+x_{3}^{2}\big)
     +x_{2}\big(-x_{1}^{3}-x_{1}x_{2}^{2}\big)\\
   &=x_{1}^{3}x_{2}+x_{1}x_{3}^{2}-x_{1}^{3}x_{2}-x_{1}x_{2}^{3}
    =x_{1}\big(x_{3}^{2}-x_{2}^{3}\big),
\end{aligned}
\]
the $x_{1}^{3}x_{2}$ terms cancelling. Hence
\[
   \boxed{\ \dot\theta=-\frac{A\sin t}{r}
   -\frac{A\cos t}{r^{3}}\,x_{1}\big(x_{3}^{2}-x_{2}^{3}\big),
   \qquad r=\sqrt{x_{1}^{2}+x_{2}^{2}}.\ }
\]
On the axis $x_{1}=x_{2}=0$ both terms are singular, as $\theta$ itself
is; off the axis the first term is the cooling seen by a fixed
thermometer and the second is what the moving material point picks up by
travelling through the spatial gradient.
\end{solution}

%-----------------------------------------------------------------------
\begin{problem}[The convective acceleration in Lamb form]
Show $\dot\bv=\bv'+\tfrac12\grad\big(|\bv|^{2}\big)+2\,\bw\times\bv$.
\end{problem}

\begin{solution}
The strategy is to compute both sides in terms of $\bD$ and $\bW$ and see
that they agree. The left side comes straight from Chapter 2.

\emph{The left side.} Apply \eqr{2.20} to the vector field $\bv$ itself,
which is \eqr{2.21}, and then split $\bL$ into its symmetric and skew
parts by \eqr{4.7}:
\[
   \dot\bv=\bv'+(\grad\bv)\bv=\bv'+\bL\bv=\bv'+\bD\bv+\bW\bv .
\]

\emph{The gradient term.} Compute $\tfrac12\grad|\bv|^{2}$ in components.
Write $|\bv|^{2}=v_{k}v_{k}$ and differentiate with the product rule,
noting that the two terms produced are equal because $k$ is summed:
\[
   \Big[\tfrac12\grad|\bv|^{2}\Big]_{i}
   =\tfrac12\big(v_{k}v_{k}\big)_{,i}
   =\tfrac12\big(v_{k,i}v_{k}+v_{k}v_{k,i}\big)
   =v_{k}v_{k,i}.
\]
Now recognise what $v_{k}v_{k,i}$ is. Since $L_{ki}=v_{k,i}$, the
expression is $L_{ki}v_{k}=(\bL^{t})_{ik}v_{k}=(\bL^{t}\bv)_{i}$, the
transpose appearing because the summed index sits in the \emph{first} slot
of $\bL$. So
\[
   \tfrac12\grad|\bv|^{2}=\bL^{t}\bv=(\bD+\bW)^{t}\bv=(\bD-\bW)\bv ,
\]
using $\bD^{t}=\bD$ and $\bW^{t}=-\bW$.

\emph{The cross-product term.} By the definition of the axial vector,
Proposition~\ref{prop:axial} with $\bW$ in place of $\bA$, we have
$\bW\bv=\bw\times\bv$ for every $\bv$.

\emph{Putting them together.} Substitute both into the claimed right side:
\[
   \bv'+\tfrac12\grad|\bv|^{2}+2\,\bw\times\bv
   =\bv'+(\bD-\bW)\bv+2\bW\bv
   =\bv'+\bD\bv+\bW\bv ,
\]
since $-\bW\bv+2\bW\bv=\bW\bv$. This is exactly the expression computed
for $\dot\bv$ above, so the two sides agree and the identity is proved.

\emph{Why anyone bothers.} The convective term $(\grad\bv)\bv$ is
quadratic in the velocity and is the source of every difficulty in fluid
mechanics. The Lamb form splits it into two pieces of quite different
character. The first, $\tfrac12\grad|\bv|^{2}$, is a gradient, so it
integrates to a potential and contributes nothing to the circulation
around a closed material curve, exactly as was seen in Problem 1 of
Chapter 5. The second, $2\bw\times\bv$, is perpendicular to $\bv$ and
therefore does no work on the flow; it is the term that vanishes when the
motion is irrotational. In a steady irrotational flow both simplifications
apply at once, $\bv'=\bze$ and $\bw=\bze$, so $\dot\bv$ is a pure gradient
and the equation of motion of Chapter 6 integrates immediately. That
integral is Bernoulli's theorem, and this rearrangement is where it comes
from.
\end{solution}

%-----------------------------------------------------------------------
\begin{problem}[A given velocity field]
For $v_{1}=ax_{2}x_{3}$, $v_{2}=-ax_{1}x_{3}$, $v_{3}=bx_{3}$ find
(a) $\bL$, (b) $\bD,\bW$, (c) $\bw$, (d) the conditions on $a,b$ for the
motion to be (i) isochoric and (ii) irrotational.
\end{problem}

\begin{solution}
\textbf{(a)} $L_{ij}=v_{i,j}$, differentiating each component with respect
to each coordinate:
\[
   (L_{ij})=\begin{pmatrix}
      0&ax_{3}&ax_{2}\\
     -ax_{3}&0&-ax_{1}\\
      0&0&b\end{pmatrix}.
\]

\textbf{(b)} $\bD=\tfrac12(\bL+\bL^{t})$ and $\bW=\tfrac12(\bL-\bL^{t})$,
entry by entry. The $(1,2)$ entries cancel in $\bD$ and double in $\bW$:
\[
   (D_{ij})=\begin{pmatrix}
      0&0&\tfrac12 ax_{2}\\
      0&0&-\tfrac12 ax_{1}\\
      \tfrac12 ax_{2}&-\tfrac12 ax_{1}&b\end{pmatrix},
   \qquad
   (W_{ij})=\begin{pmatrix}
      0&ax_{3}&\tfrac12 ax_{2}\\
     -ax_{3}&0&-\tfrac12 ax_{1}\\
     -\tfrac12 ax_{2}&\tfrac12 ax_{1}&0\end{pmatrix}.
\]

\textbf{(c)} From $\bW\bu=\bw\times\bu$ one has $W_{ij}=-\eps_{ijk}w_{k}$,
so $w_{1}=-W_{23}$, $w_{2}=W_{13}$, $w_{3}=-W_{12}$:
\[
   \bw=\tfrac{a}{2}\big(x_{1}\be_{1}+x_{2}\be_{2}\big)-a x_{3}\be_{3}.
\]
Check against \eqr{4.23}, $\bw=\tfrac12\crl\bv$:
\[
\begin{aligned}
   (\crl\bv)_{1}&=v_{3,2}-v_{2,3}=0+ax_{1}=ax_{1},\\
   (\crl\bv)_{2}&=v_{1,3}-v_{3,1}=ax_{2}-0=ax_{2},\\
   (\crl\bv)_{3}&=v_{2,1}-v_{1,2}=-ax_{3}-ax_{3}=-2ax_{3},
\end{aligned}
\]
whose half is the $\bw$ just found.

\textbf{(d)} (i) Isochoric means $\dot J=0$, i.e.\ $\tr\bL=0$ by
Problem 2. Here $\tr\bL=0+0+b=b$, so the condition is
\[
   b=0 .
\]
(ii) Irrotational means $\bw=\bze$ at every point; the component
$w_{1}=\tfrac12 ax_{1}$ must vanish for all $x_{1}$, so
\[
   a=0 ,
\]
and then $\bw=\bze$ whatever $b$ is. Note the two conditions are
independent: $b$ controls the volume change, $a$ the rotation.
\end{solution}

%-----------------------------------------------------------------------
\begin{problem}[Spin of a fibre when $\bW=\Ob$]
Let $\bF=\lambda\,\be_{1}\otimes\hbE_{1}
+\lambda^{-1/2}\big(\be_{2}\otimes\hbE_{2}+\be_{3}\otimes\hbE_{3}\big)$
with $\lambda=\lambda(t)$.
(a) Show the motion is isochoric. (b) Verify the stated $\bF^{-1}$.
(c) Find $\bL$ and show $\bW=\Ob$. (d) With
$\bm=\cos\phi\,\ber(\theta)+\sin\phi\,\be_{3}$, find $\dot\theta$ and
$\dot\phi$, and identify the directions with $\dot\bm=\bze$.
\end{problem}

\begin{solution}
\textbf{(a)} $\bF$ is already in the form \eqr{3.43} with orthonormal frames
on both sides, so its principal stretches are the coefficients and
\[
   J=\lambda\cdot\lambda^{-1/2}\cdot\lambda^{-1/2}
   =\lambda^{1-\frac12-\frac12}=1 .
\]

\textbf{(b)} Multiply, using $(\ba\otimes\bb)(\bc\otimes\bd)
=(\ip\bb\bc)\ba\otimes\bd$; only the three diagonal products survive:
\[
   \bF\bF^{-1}=\lambda\lambda^{-1}\be_{1}\otimes\be_{1}
   +\lambda^{-1/2}\lambda^{1/2}\big(\be_{2}\otimes\be_{2}
   +\be_{3}\otimes\be_{3}\big)=\I .
\]

\textbf{(c)} Differentiate $\bF$ at fixed $\bX$, the bases being constant:
\[
   \dot\bF=\dot\lambda\,\be_{1}\otimes\hbE_{1}
   -\tfrac12\lambda^{-3/2}\dot\lambda
   \big(\be_{2}\otimes\hbE_{2}+\be_{3}\otimes\hbE_{3}\big),
\]
and multiply by $\bF^{-1}$ of part (b):
\[
   \boxed{\ \bL=\dot\bF\bF^{-1}
   =\frac{\dot\lambda}{\lambda}\,\be_{1}\otimes\be_{1}
   -\frac{\dot\lambda}{2\lambda}
   \big(\be_{2}\otimes\be_{2}+\be_{3}\otimes\be_{3}\big).\ }
\]
This is symmetric, so $\bW=\Skw\bL=\Ob$ and $\bD=\bL$. Its trace is
$\dot\lambda/\lambda-\dot\lambda/\lambda=0$, consistent with (a) and
Problem 2.

\textbf{(d)} Write $\alpha=\dot\lambda/\lambda$, so
$\bD=\alpha\be_{1}\otimes\be_{1}
-\tfrac\alpha2(\be_{2}\otimes\be_{2}+\be_{3}\otimes\be_{3})$, and
\[
\begin{aligned}
   \bm&=\big(\cos\phi\cos\theta,\ \cos\phi\sin\theta,\ \sin\phi\big),\\
   \bD\bm&=\big(\alpha\cos\phi\cos\theta,\
   -\tfrac\alpha2\cos\phi\sin\theta,\ -\tfrac\alpha2\sin\phi\big).
\end{aligned}
\]
Differentiating the parametrisation, with $\beth=\ber'(\theta)$ and
$\beph=-\sin\phi\,\ber+\cos\phi\,\be_{3}$,
\[
   \dot\bm=\cos\phi\,\dot\theta\,\beth+\dot\phi\,\beph
   \ \Longrightarrow\
   \dot\theta=\frac{\ip{\beth}{\dot\bm}}{\cos\phi},
   \qquad
   \dot\phi=\ip{\beph}{\dot\bm},
\]
since $\{\bm,\beth,\beph\}$ is orthonormal. By \eqr{4.16} with $\bW=\Ob$,
$\dot\bm=\bD\bm-(\ip\bm{\bD\bm})\bm$, and both $\beth$ and $\beph$ are
perpendicular to $\bm$, so the second term drops out of both projections:
$\ip{\beth}{\dot\bm}=\ip{\beth}{\bD\bm}$ and
$\ip{\beph}{\dot\bm}=\ip{\beph}{\bD\bm}$. Now
\[
\begin{aligned}
   \ip{\beth}{\bD\bm}
   &=-\sin\theta\,\alpha\cos\phi\cos\theta
     +\cos\theta\big(-\tfrac\alpha2\cos\phi\sin\theta\big)\\
   &=-\tfrac{3\alpha}{2}\cos\phi\sin\theta\cos\theta
    =-\tfrac{3\alpha}{4}\cos\phi\,\sin2\theta,
\end{aligned}
\]
\[
\begin{aligned}
   \ip{\beph}{\bD\bm}
   &=-\alpha\sin\phi\cos\phi\cos^{2}\theta
     +\tfrac\alpha2\sin\phi\cos\phi\sin^{2}\theta
     -\tfrac\alpha2\sin\phi\cos\phi\\
   &=\alpha\sin\phi\cos\phi
     \Big[-\cos^{2}\theta+\tfrac12\sin^{2}\theta-\tfrac12\Big]
    =-\tfrac{3\alpha}{2}\sin\phi\cos\phi\,\cos^{2}\theta,
\end{aligned}
\]
using $\tfrac12\sin^{2}\theta-\tfrac12=-\tfrac12\cos^{2}\theta$. Hence
\[
   \boxed{\ \dot\theta=-\frac{3\dot\lambda}{4\lambda}\,\sin2\theta,
   \qquad
   \dot\phi=-\frac{3\dot\lambda}{4\lambda}\,\sin2\phi\,\cos^{2}\theta .\ }
\]
Both are generally non-zero although $\bW=\Ob$: this is
Remark~\ref{rem:Wtrap} made quantitative, and it is the point of the
problem.

\emph{When is $\dot\bm=\bze$?} Since $\bW=\Ob$, \eqr{4.16} gives
$\dot\bm=\bD\bm-(\ip\bm{\bD\bm})\bm$, which vanishes exactly when
$\bD\bm$ is parallel to $\bm$, i.e.\ when $\bm$ is a principal direction
of $\bD$. Reading the boxed formulas, $\dot\theta=0$ needs
$\sin2\theta=0$ and $\dot\phi=0$ needs $\sin2\phi\cos^{2}\theta=0$, so
\[
\begin{aligned}
   \theta&=\tfrac\pi2\ \text{or}\ \tfrac{3\pi}{2}:
     &&\cos\theta=0,\ \text{so any }\phi\ \text{works},\\
   \theta&=0\ \text{or}\ \pi: &&\phi=0\ \text{or}\ \tfrac\pi2 .
\end{aligned}
\]
The first family is every direction in the $(\be_{2},\be_{3})$ plane, the
second gives $\bm=\pm\be_{1}$ and $\bm=\pm\be_{3}$. That is exactly the
principal set of $\bD$: the value $-\alpha/2$ is repeated, so the whole
$(\be_{2},\be_{3})$ plane is principal, and $\be_{1}$ carries the isolated
value $\alpha$.
\end{solution}

%-----------------------------------------------------------------------
\begin{problem}[Rate of area change and rotation of the normal]
With $\alpha=\dl a/\dl A$ the local area ratio of a material surface,
(a) show $\alpha>0$; (b) establish
$\dot\alpha/\alpha=\tr\bD-\ip\bn{\bD\bn}$ and
$\dot\bn=(\ip\bn{\bD\bn})\bn-\bL^{t}\bn$.
\end{problem}

\begin{solution}
\textbf{(a)} By Piola--Nanson \eqr{2.54}, $\alpha\bn=\bF^{*}\bN$ with
$|\bn|=1$, so $\alpha=|\bF^{*}\bN|$. Now $\bF^{*}=J\bF^{-t}$ is
invertible, since $\det\bF^{*}=J^{2}\neq0$, and $\bN\neq\bze$; hence
$\bF^{*}\bN\neq\bze$ and $\alpha>0$.

\textbf{(b)} Differentiate $\alpha\bn=J\bF^{-t}\bN$ at fixed $\bN$, which
is legitimate because $\bN$ is a material vector. Two ingredients are
needed. First $\dot J=J\tr\bD$ from Problem 2. Second, by Problem 1(b),(d)
and $\dot\bF=\bL\bF$,
\[
   \big(\bF^{-t}\big)^{\textstyle\cdot}
   =\Big(\big(\bF^{-1}\big)^{\textstyle\cdot}\Big)^{t}
   =\big(-\bF^{-1}\bL\bF\bF^{-1}\big)^{t}
   =-\big(\bF^{-1}\bL\big)^{t}
   =-\bL^{t}\bF^{-t}.
\]
Therefore
\[
   \dot\alpha\,\bn+\alpha\dot\bn
   =\dot J\bF^{-t}\bN+J\big(\bF^{-t}\big)^{\textstyle\cdot}\bN
   =(\tr\bD)\,\alpha\bn-\bL^{t}\alpha\bn ,
\]
having recognised $J\bF^{-t}\bN=\alpha\bn$ twice. Dot with $\bn$, using
$\ip{\bn}{\dot\bn}=0$ and
$\ip{\bn}{\bL^{t}\bn}=\ip{\bn}{\bL\bn}=\ip\bn{\bD\bn}$:
\[
   \boxed{\ \frac{\dot\alpha}{\alpha}=\tr\bD-\ip\bn{\bD\bn}\ }
\]
Substituting this back and dividing by $\alpha$,
\[
   \dot\bn=(\tr\bD)\bn-\bL^{t}\bn-\big(\tr\bD-\ip\bn{\bD\bn}\big)\bn
   =\boxed{\ (\ip\bn{\bD\bn})\bn-\bL^{t}\bn .\ }
\]
\begin{figure}[htbp]\centering
\renewcommand{\thefigure}{E}
\begin{tikzpicture}[scale=1.5,
   ar/.style={-{Stealth[length=2.4mm]},thick}]
% the material square now, and a moment later under simple shear
\draw[guide] (0,0) rectangle (1.5,1.5);
\draw[thick] (0,0)--(1.5,0)--(2.05,1.5)--(0.55,1.5)--cycle;
% the left edge as a material fibre
\draw[ar] (0,0)--(0.55,1.5);
\node[font=\small,anchor=south east] at (0.56,1.42) {$\bm$};
\draw[guide] (0,0)--(0,1.5);
% its normal, before and after
\draw[ar] (0.275,0.75)--(1.20,0.411);
\node[font=\small,anchor=north west=-1pt] at (1.17,0.40) {$\bn$};
\draw[guide] (0.275,0.75)--(1.20,0.75);
% right angle between m and n
\draw[gray!65,thin] (0.363,0.99)--(0.487,1.048)--(0.545,0.924);
\node[font=\footnotesize,text=gray!80,anchor=north] at (1.0,-0.22)
   {dashed horizontal: where convecting $\bn$ with $\bL$ would leave it};
\node[font=\small,anchor=north] at (1.0,-0.68)
   {$\dot\bm=\bL\bm-(\bm\cdot\bD\bm)\bm$
    \qquad
    $\dot\bn=(\bn\cdot\bD\bn)\bn-\bL^{t}\bn$};
\end{tikzpicture}
\caption{Simple shear, with the square as it stands now dashed. The left
edge is a material fibre and also a material surface; its unit tangent
$\bm$ and unit normal $\bn$ must stay perpendicular, and they do. That is
what the transpose in $-\bL^{t}\bn$ buys: here
$\bL=\kappa\,\be_{1}\otimes\be_{2}$, so $\bL\bm$ tips the fibre to the
right while $-\bL^{t}\bn=-\kappa\be_{2}$ tips the normal down by the same
angle. Convecting the normal with $\bL$ instead, the common slip, gives
$\bL\be_{1}=\bze$: the normal would not move at all, and would stop being
normal.}
\label{fig:linevsnormal}
\end{figure}

Compare with \eqr{4.16}: a material \emph{line} element is convected by
$\bL$, a material \emph{normal} by $-\bL^{t}$. That sign and transpose
are the rate form of the difference between $\bF$ and $\bF^{-t}$ in
Piola--Nanson, and forgetting them is the standard error.
\end{solution}

%-----------------------------------------------------------------------
\begin{problem}[Stretching when the principal axes are frozen in the body]
If $\dot\bu_{i}=\bze$, show
$\bD=\sum_{i}(\ln\lambda_{i})^{\textstyle\cdot}\,\bv_{i}\otimes\bv_{i}$.
\end{problem}

\begin{solution}
With $\bU=\sum_{i}\lambda_{i}\bu_{i}\otimes\bu_{i}$ and
$\dot\bu_{i}=\bze$, only the coefficients move:
\[
   \dot\bU=\sum_{i}\dot\lambda_{i}\,\bu_{i}\otimes\bu_{i},
   \qquad
   \bU^{-1}=\sum_{i}\lambda_{i}^{-1}\,\bu_{i}\otimes\bu_{i},
\]
the inverse by Remark~\ref{rem:powers}. Now multiply them. The product of
the two sums has nine terms, one for each pair $(i,j)$:
\[
   \dot\bU\bU^{-1}
   =\sum_{i}\sum_{j}\frac{\dot\lambda_{i}}{\lambda_{j}}
     \big(\bu_{i}\otimes\bu_{i}\big)\big(\bu_{j}\otimes\bu_{j}\big).
\]
Each product of dyads collapses by the rule
$(\ba\otimes\bb)(\bc\otimes\bd)=(\ip\bb\bc)\,\ba\otimes\bd$, which is the
definition of the tensor product applied twice. Here $\bb=\bu_{i}$ and
$\bc=\bu_{j}$, so the scalar produced is $\ip{\bu_{i}}{\bu_{j}}=\dd_{ij}$
by orthonormality of the principal basis. That kills every term with
$i\neq j$, six of the nine, and leaves the three diagonal ones with
$\lambda_{j}$ replaced by $\lambda_{i}$:
\[
   \dot\bU\bU^{-1}
   =\sum_{i}\frac{\dot\lambda_{i}}{\lambda_{i}}\,
     \bu_{i}\otimes\bu_{i}
   =\sum_{i}(\ln\lambda_{i})^{\textstyle\cdot}\,\bu_{i}\otimes\bu_{i},
\]
the last step being the same chain rule that produced \eqr{4.15},
$\dot\lambda/\lambda=(\ln\lambda)^{\textstyle\cdot}$. The result is
\emph{symmetric}, being a sum of terms $\bu\otimes\bu$ with real
coefficients, and that observation is what the rest of the solution turns
on. Now look at \eqr{4.5}:
\[
   \bL=\underbrace{\dot\bR\bR^{t}}_{\text{skew, Problem 1(c)}}
      +\underbrace{\bR\big(\dot\bU\bU^{-1}\big)\bR^{t}}_{\text{symmetric}},
\]
the second being symmetric because $\bR\bS\bR^{t}$ is symmetric whenever
$\bS$ is. By the uniqueness of the symmetric--skew splitting, Problem 14
of Chapter 1, these two \emph{are} $\bD$ and $\bW$. Hence, using
$(\bR\ba)\otimes(\bR\bb)=\bR(\ba\otimes\bb)\bR^{t}$ from \eqr{3.41} and
$\bv_{i}=\bR\bu_{i}$,
\[
   \bD=\bR\Big(\sum_{i}(\ln\lambda_{i})^{\textstyle\cdot}
       \bu_{i}\otimes\bu_{i}\Big)\bR^{t}
   =\boxed{\ \sum_{i}(\ln\lambda_{i})^{\textstyle\cdot}\,
     \bv_{i}\otimes\bv_{i}\ }
\]
and $\bW=\dot\bR\bR^{t}$.

Two readings. The principal values of $\bD$ are the logarithmic stretch
rates, in agreement with \eqr{4.17}; and $\bD$ is here coaxial with $\bV$,
which is exactly the case in which the Hencky strain of
Remark~\ref{rem:sethhill} satisfies $\bD=(\ln\bV)^{\textstyle\cdot}$. In
general, when $\dot\bu_{i}\neq\bze$, that last identity is false.
\end{solution}

%-----------------------------------------------------------------------
\begin{problem}[An azimuthally graded axial flow]
For $\bv=(\theta/\theta_{0})W\be_{3}$ with $W,\theta_{0}$ constant, find
(a) $\bL,\bD,\bW$; (b) the principal values and vectors of $\bD$, and
$\bw$.
\end{problem}

\begin{solution}
\textbf{(a)} $\bv=f\be_{3}$ with $f=W\theta/\theta_{0}$ a scalar field and
$\be_{3}$ constant, so by Problem 24(f) of Chapter 1,
$\grad(f\be_{3})=\be_{3}\otimes\grad f$. By Problem 25 of Chapter 1 in
cylindrical form, $\grad\theta=\bet/r$, whence
\[
   \boxed{\ \bL=\frac{W}{\theta_{0}r}\,\be_{3}\otimes\bet\ }
\]
and, splitting by \eqr{4.7} with $c:=W/(2\theta_{0}r)$,
\[
   \bD=c\big(\be_{3}\otimes\bet+\bet\otimes\be_{3}\big),
   \qquad
   \bW=c\big(\be_{3}\otimes\bet-\bet\otimes\be_{3}\big).
\]
Note $\tr\bD=2c\,\ip{\be_{3}}{\bet}=0$: the flow is isochoric.

\textbf{(b)} In the ordered orthonormal frame $(\ber,\bet,\be_{3})$,
\[
   (D)=\begin{pmatrix}0&0&0\\ 0&0&c\\ 0&c&0\end{pmatrix},
\]
a pure shear on the $(\bet,\be_{3})$ pair with $\ber$ inert. Its
characteristic equation is $-\eta(\eta^{2}-c^{2})=0$, so
\[
   \eta=0,\ +c,\ -c;
   \qquad
   \bnu_{0}=\ber,\quad
   \bnu_{\pm}=\tfrac{\sqrt2}{2}\big(\bet\pm\be_{3}\big).
\]
For the vorticity, $\bW=\ba\otimes\bb-\bb\otimes\ba$ has axial vector
$\bb\times\ba$, since
\[
   (\ba\otimes\bb-\bb\otimes\ba)\bu=(\ip\bb\bu)\ba-(\ip\ba\bu)\bb
   =(\bb\times\ba)\times\bu
\]
by Problem 11(a). With $\ba=c\,\be_{3}$, $\bb=\bet$ and
$\bet\times\be_{3}=\ber$,
\[
   \boxed{\ \bw=\frac{W}{2\theta_{0}r}\,\ber\ }
\]
which checks against \eqr{4.23}: $\crl(f\be_{3})=\grad f\times\be_{3}
=\tfrac{W}{\theta_{0}r}\bet\times\be_{3}=\tfrac{W}{\theta_{0}r}\ber$, half
of which is $\bw$.
\end{solution}

\begin{remark}[the field of Problem 9 is not single valued]
\label{rem:P9multi}
$\theta$ jumps by $2\pi$ on crossing the half plane $\theta=0$, so
$\bv=(W\theta/\theta_{0})\be_{3}$ jumps by $2\pi W/\theta_{0}$ there: the
velocity field is discontinuous across that surface, and the motion it
describes tears the material. The \emph{gradient} $\bL$ is nevertheless
perfectly single valued, because $\grad\theta=\bet/r$ is, which is why the
computation above goes through without complaint. This is a useful warning:
$\bL$, $\bD$, $\bW$ can be smooth where $\bv$ is not, and a velocity field
must be checked for single-valuedness separately.
\end{remark}

%-----------------------------------------------------------------------
\begin{problem}[Circular flow]
For $\bv=rf(r)\bet$ in cylindrical coordinates, show
$\bL=\big[\partial_{r}(rf)\big]\bet\otimes\ber-f\,\ber\otimes\bet$ and
find $\bD$ and $\bW$.
\end{problem}

\begin{solution}
Apply \eqr{1.145}, the polar form of $\grad\bu$, to
$u_{r}=0$, $u_{\theta}=rf(r)$, $u_{z}=0$, all independent of $\theta$ and
$z$. The six coefficients are
\[
\begin{aligned}
   u_{r,r}&=0, &\tfrac1r\big(u_{r,\theta}-u_{\theta}\big)&=-f,
   &u_{\theta,r}&=(rf)',\\
   \tfrac1r\big(u_{\theta,\theta}+u_{r}\big)&=0,
   &u_{z,r}&=0, &\tfrac1r u_{z,\theta}&=0,
\end{aligned}
\]
so
\[
   \boxed{\ \bL=(rf)'\,\bet\otimes\ber-f\,\ber\otimes\bet\ }
\]
as stated. Splitting by \eqr{4.7}, and using $(rf)'=f+rf'$ so that
$(rf)'-f=rf'$ and $(rf)'+f=2f+rf'$,
\[
   \bD=\frac{rf'}{2}\big(\ber\otimes\bet+\bet\otimes\ber\big),
   \qquad
   \bW=\frac{2f+rf'}{2}\big(\bet\otimes\ber-\ber\otimes\bet\big).
\]
By the axial-vector rule of Problem 9, with $\ba=\bet$, $\bb=\ber$ and
$\ber\times\bet=\be_{3}$,
\[
   \bw=\frac{2f+rf'}{2}\,\be_{3}=\frac{1}{2r}\big(r^{2}f\big)'\,\be_{3}.
\]

Two special cases show what the two tensors separate. If $f$ is constant
the flow is a rigid rotation: then $f'=0$, so $\bD=\Ob$, as
Remark~\ref{rem:Drigid} requires, and $\bw=f\be_{3}$, the angular
velocity. If instead $f=k/r^{2}$, then $r^{2}f=k$ is constant and
$\bw=\bze$: the flow is irrotational although every particle travels in a
circle. This is the potential vortex, and it is the standard
counterexample to the belief that curved streamlines imply vorticity.
Its stretching is $\bD=-\tfrac{k}{r^{2}}(\ber\otimes\bet
+\bet\otimes\ber)$, which is not zero: the material is being sheared even
though it is not spinning.

The two cases are worth seeing side by side, because they separate the two
things $\bL$ measures. Put a radial and an azimuthal material fibre at the
same point and let $\Omega$ denote the rate at which each turns about
$\be_{3}$. From $\dot\bm=\bL\bm-(\bm\cdot\bD\bm)\bm$ with
$\ber\cdot\bD\ber=\bet\cdot\bD\bet=0$,
\[
   \dot\ber=(rf)'\,\bet,\qquad \dot\bet=-f\,\ber,
   \qquad\text{so}\qquad
   \Omega_{r}=(rf)'=f+rf',\quad \Omega_{\theta}=f .
\]
Their mean is the vorticity and their difference is the shear:
\[
   \tfrac12\big(\Omega_{r}+\Omega_{\theta}\big)=\frac{2f+rf'}{2},
   \qquad
   \Omega_{\theta}-\Omega_{r}=-rf'=-2\,\ber\cdot\bD\bet ,
\]
in agreement with $\bw$ above and with \eqr{4.21}. Rigid rotation is
$\Omega_{r}=\Omega_{\theta}=f$; the potential vortex $f=k/r^{2}$ has
$rf'=-2f$, hence $\Omega_{r}=-f$ and $\Omega_{\theta}=+f$, equal and
opposite.
\end{solution}

\begin{figure}[htbp]\centering
\renewcommand{\thefigure}{F}
\begin{tikzpicture}[scale=1.0,
   ar/.style={-{Stealth[length=2.3mm]},thick},
   old/.style={-{Stealth[length=2.0mm]},semithick,gray!55,densely dashed},
   sw/.style={-{Stealth[length=1.9mm]},thin,gray!80}]
% ================= panel 1 : rigid rotation =================
\begin{scope}[xshift=0cm]
  % a short arc of the streamline through P; its centre lies 2.2 to the
  % left, so at P the outward radial direction is +x and the azimuthal +y
  \draw[guide,-{Stealth[length=2.0mm]}] (-0.295,-1.10) arc (-30:30:2.2);
  \fill (0,0) circle (1.3pt);
  \node[font=\small,anchor=north east] at (-0.05,-0.07) {$P$};
  % the two fibres as they stand now
  \draw[old] (0,0)--(1.25,0);  \draw[old] (0,0)--(0,1.25);
  \node[font=\small,text=gray!65,below] at (1.16,-0.07) {$\ber$};
  \node[font=\small,text=gray!65,right] at (0.07,1.14) {$\bet$};
  % both turn by the same +25 deg
  \draw[ar] (0,0)--(25:1.25);
  \draw[ar] (0,0)--(115:1.25);
  \draw[sw] (1.44,0) arc (0:25:1.44);
  \draw[sw] (0,1.44) arc (90:115:1.44);
  \draw[gray!65,thin] (25:0.34) arc (25:115:0.34);
  \node[font=\small,anchor=north] at (0.45,-1.42)
     {$\Omega_{r}=\Omega_{\theta}=f$};
  \node[font=\small,anchor=north] at (0.45,-1.95)
     {$\bD=\Ob$,\quad $\bw=f\be_{3}$};
  \node[font=\footnotesize,text=gray!80,align=center,anchor=north]
     at (0.45,-2.48) {rigid rotation:\\ it spins, it does not shear};
\end{scope}
% ================= panel 2 : potential vortex =================
\begin{scope}[xshift=5.9cm]
  \draw[guide,-{Stealth[length=2.0mm]}] (-0.295,-1.10) arc (-30:30:2.2);
  \fill (0,0) circle (1.3pt);
  \node[font=\small,anchor=north east] at (-0.05,-0.07) {$P$};
  \draw[old] (0,0)--(1.25,0);  \draw[old] (0,0)--(0,1.25);
  \node[font=\small,text=gray!65,above=-1pt] at (1.16,0.04) {$\ber$};
  \node[font=\small,text=gray!65,right] at (0.07,1.14) {$\bet$};
  % er turns BACK by 25, et turns FORWARD by 25: the pair opens
  \draw[ar] (0,0)--(-25:1.25);
  \draw[ar] (0,0)--(115:1.25);
  \draw[sw] (1.44,0) arc (0:-25:1.44);
  \draw[sw] (0,1.44) arc (90:115:1.44);
  \draw[gray!65,thin] (-25:0.34) arc (-25:115:0.34);
  \node[font=\small,anchor=north] at (0.45,-1.42)
     {$\Omega_{r}=-f$,\quad $\Omega_{\theta}=+f$};
  \node[font=\small,anchor=north] at (0.45,-1.95)
     {$\bw=\bze$,\quad $\bD\neq\Ob$};
  \node[font=\footnotesize,text=gray!80,align=center,anchor=north]
     at (0.45,-2.48) {potential vortex:\\ it shears, it does not spin};
\end{scope}
\end{tikzpicture}
\caption{Problem 10, at one point $P$ of a circular streamline, with the
two material fibres as they stand now shown dashed. In rigid rotation both
turn at the same rate, so their mean rate, the vorticity, is $f$ and their
relative rate, the shearing, is zero. In the potential vortex $f=k/r^{2}$
the two rates are equal and opposite: the mean vanishes, so a paddle wheel
placed at $P$ would circle the origin without ever turning, yet the right
angle between the fibres opens at rate $2f$, so the material is being
sheared throughout. Curved streamlines therefore say nothing about
vorticity.}
\label{fig:vortex}
\end{figure}
